%%%%%%%%%%%%%%%%%%%%%%%%%%%%%%%%%%%%%%%%%%%%%%%%%%%%%%%%%%%%
% 4.2 ALGEBRAIC TOPOLOGY
% The Composition of Advanced Mathematics
%%%%%%%%%%%%%%%%%%%%%%%%%%%%%%%%%%%%%%%%%%%%%%%%%%%%%%%%%%%%

\section{Algebraic Topology}
\label{sec:algebraic-topology}

\prerequisite{
Point-Set Topology, Group Theory, basic Ring Theory,
Linear Algebra, Mathematical Logic, and Proof Techniques.
}

\learningobjectives{
By the end of this section, the reader should be able to:
\begin{itemize}
    \item explain the central philosophy of algebraic topology;
    \item define paths, loops, and homotopies;
    \item distinguish homotopy from homeomorphism;
    \item define homotopy equivalence and contractibility;
    \item explain deformation retracts;
    \item define the fundamental group;
    \item compute fundamental groups of several elementary spaces;
    \item explain induced homomorphisms on fundamental groups;
    \item describe covering spaces and lifting properties;
    \item state the relationship between covering spaces and fundamental groups;
    \item explain simply connected spaces and universal covers;
    \item define simplices and simplicial complexes;
    \item construct simplicial chain groups;
    \item define boundary homomorphisms;
    \item explain cycles and boundaries;
    \item define simplicial homology groups;
    \item compute basic homology groups;
    \item interpret Betti numbers geometrically;
    \item explain reduced homology;
    \item describe singular homology conceptually;
    \item state the Euler characteristic in terms of cells and homology;
    \item explain the role of exact sequences;
    \item state the Mayer--Vietoris principle;
    \item describe cohomology and the cup product;
    \item recognize homotopy and homology as topological invariants;
    \item explain how algebraic topology connects geometry with algebra.
\end{itemize}
}

%%%%%%%%%%%%%%%%%%%%%%%%%%%%%%%%%%%%%%%%%%%%%%%%%%%%%%%%%%%%
\subsection{The Central Idea of Algebraic Topology}
\label{subsec:central-idea-algebraic-topology}
%%%%%%%%%%%%%%%%%%%%%%%%%%%%%%%%%%%%%%%%%%%%%%%%%%%%%%%%%%%%

Point-set topology provides a language for discussing spaces using

\[
\boxed{
\text{open sets},
\quad
\text{continuity},
\quad
\text{compactness},
\quad
\text{connectedness}.
}
\]

Algebraic topology goes further.

It assigns algebraic objects to topological spaces.

Schematically,

\[
\boxed{
\text{topological space}
\longrightarrow
\text{algebraic invariant}.
}
\]

These algebraic objects may include

\[
\boxed{
\text{groups},
\quad
\text{abelian groups},
\quad
\text{rings},
\quad
\text{modules}.
}
\]

The resulting algebra can reveal properties of a space that may be difficult
to detect directly.

\begin{conceptbox}
The basic philosophy of algebraic topology is:

\[
\boxed{
\text{translate topology into algebra}.
}
\]

If two spaces produce different algebraic invariants, then the spaces cannot
be topologically equivalent.
\end{conceptbox}

%%%%%%%%%%%%%%%%%%%%%%%%%%%%%%%%%%%%%%%%%%%%%%%%%%%%%%%%%%%%
\subsection{Topological Invariants}
\label{subsec:algebraic-topological-invariants}
%%%%%%%%%%%%%%%%%%%%%%%%%%%%%%%%%%%%%%%%%%%%%%%%%%%%%%%%%%%%

A \emph{topological invariant} is a quantity or algebraic structure preserved
by homeomorphisms.

Examples developed in algebraic topology include

\[
\boxed{
\pi_1(X),
\qquad
H_n(X),
\qquad
H^n(X).
}
\]

Here:

\[
\pi_1(X)
\]

denotes the fundamental group,

\[
H_n(X)
\]

denotes the \(n\)-th homology group, and

\[
H^n(X)
\]

denotes the \(n\)-th cohomology group.

The lowercase Greek letter

\[
\pi
\]

is read ``pi.''

%%%%%%%%%%%%%%%%%%%%%%%%%%%%%%%%%%%%%%%%%%%%%%%%%%%%%%%%%%%%
\subsection{Paths}
\label{subsec:algebraic-topology-paths}
%%%%%%%%%%%%%%%%%%%%%%%%%%%%%%%%%%%%%%%%%%%%%%%%%%%%%%%%%%%%

Recall that a path in a topological space \(X\) is a continuous map

\[
\boxed{
\gamma:[0,1]\to X.
}
\]

The endpoints are

\[
\gamma(0)
\qquad\text{and}\qquad
\gamma(1).
\]

If

\[
\gamma(0)=x
\qquad\text{and}\qquad
\gamma(1)=y,
\]

then \(\gamma\) is called a path from \(x\) to \(y\).

%%%%%%%%%%%%%%%%%%%%%%%%%%%%%%%%%%%%%%%%%%%%%%%%%%%%%%%%%%%%
\subsection{Loops}
\label{subsec:loops}
%%%%%%%%%%%%%%%%%%%%%%%%%%%%%%%%%%%%%%%%%%%%%%%%%%%%%%%%%%%%

\begin{definitionbox}{Loop}
A \emph{loop based at \(x_0\)} is a continuous map

\[
\boxed{
\gamma:[0,1]\to X
}
\]

such that

\[
\boxed{
\gamma(0)=\gamma(1)=x_0.
}
\]
\end{definitionbox}

The point \(x_0\) is called the \emph{basepoint}.

Thus a loop begins and ends at the same point.

%%%%%%%%%%%%%%%%%%%%%%%%%%%%%%%%%%%%%%%%%%%%%%%%%%%%%%%%%%%%
\subsection{Constant Loops}
\label{subsec:constant-loops}
%%%%%%%%%%%%%%%%%%%%%%%%%%%%%%%%%%%%%%%%%%%%%%%%%%%%%%%%%%%%

For a point \(x_0\in X\), the constant loop is

\[
\boxed{
e_{x_0}(t)=x_0
}
\]

for every

\[
t\in[0,1].
\]

This loop plays the role of the identity element in the fundamental group.

%%%%%%%%%%%%%%%%%%%%%%%%%%%%%%%%%%%%%%%%%%%%%%%%%%%%%%%%%%%%
\subsection{Reversing a Path}
\label{subsec:reverse-path}
%%%%%%%%%%%%%%%%%%%%%%%%%%%%%%%%%%%%%%%%%%%%%%%%%%%%%%%%%%%%

If

\[
\gamma:[0,1]\to X
\]

is a path, define its reverse by

\[
\boxed{
\overline{\gamma}(t)
=
\gamma(1-t).
}
\]

The reversed path traverses the same geometric route in the opposite
direction.

If \(\gamma\) travels from \(x\) to \(y\), then

\[
\overline{\gamma}
\]

travels from \(y\) to \(x\).

%%%%%%%%%%%%%%%%%%%%%%%%%%%%%%%%%%%%%%%%%%%%%%%%%%%%%%%%%%%%
\subsection{Path Concatenation}
\label{subsec:path-concatenation}
%%%%%%%%%%%%%%%%%%%%%%%%%%%%%%%%%%%%%%%%%%%%%%%%%%%%%%%%%%%%

Suppose

\[
\alpha(1)=\beta(0).
\]

Then the concatenation

\[
\alpha*\beta
\]

is obtained by traversing \(\alpha\) first and \(\beta\) second.

One standard parametrization is

\[
\boxed{
(\alpha*\beta)(t)
=
\begin{cases}
\alpha(2t),
&
0\leq t\leq \frac12,
\\[4pt]
\beta(2t-1),
&
\frac12\leq t\leq1.
\end{cases}
}
\]

The symbol

\[
*
\]

is read ``star'' and here denotes path concatenation.

%%%%%%%%%%%%%%%%%%%%%%%%%%%%%%%%%%%%%%%%%%%%%%%%%%%%%%%%%%%%
\subsection{Homotopy of Maps}
\label{subsec:homotopy-of-maps}
%%%%%%%%%%%%%%%%%%%%%%%%%%%%%%%%%%%%%%%%%%%%%%%%%%%%%%%%%%%%

The next idea formalizes continuous deformation.

\begin{definitionbox}{Homotopy}
Let

\[
f,g:X\to Y
\]

be continuous maps.

A \emph{homotopy} from \(f\) to \(g\) is a continuous map

\[
\boxed{
H:X\times[0,1]\to Y
}
\]

such that

\[
H(x,0)=f(x)
\]

and

\[
H(x,1)=g(x)
\]

for every \(x\in X\).
\end{definitionbox}

We write

\[
\boxed{
f\simeq g.
}
\]

The symbol

\[
\simeq
\]

is read ``is homotopic to.''

%%%%%%%%%%%%%%%%%%%%%%%%%%%%%%%%%%%%%%%%%%%%%%%%%%%%%%%%%%%%
\subsection{Interpreting a Homotopy}
\label{subsec:interpreting-homotopy}
%%%%%%%%%%%%%%%%%%%%%%%%%%%%%%%%%%%%%%%%%%%%%%%%%%%%%%%%%%%%

For each fixed

\[
t\in[0,1],
\]

the function

\[
H_t(x)=H(x,t)
\]

can be regarded as an intermediate map.

Thus

\[
H_0=f
\]

and

\[
H_1=g.
\]

One may picture the parameter \(t\) as time.

Then a homotopy is a continuous family of maps transforming \(f\) into \(g\).

%%%%%%%%%%%%%%%%%%%%%%%%%%%%%%%%%%%%%%%%%%%%%%%%%%%%%%%%%%%%
\subsection{Homotopy Is an Equivalence Relation}
\label{subsec:homotopy-equivalence-relation}
%%%%%%%%%%%%%%%%%%%%%%%%%%%%%%%%%%%%%%%%%%%%%%%%%%%%%%%%%%%%

For suitable spaces and maps,

\[
\simeq
\]

is an equivalence relation.

It is:

\[
\boxed{\text{reflexive}},
\]

because \(f\simeq f\);

\[
\boxed{\text{symmetric}},
\]

because \(f\simeq g\) implies \(g\simeq f\); and

\[
\boxed{\text{transitive}},
\]

because

\[
f\simeq g
\qquad\text{and}\qquad
g\simeq h
\]

imply

\[
f\simeq h.
\]

%%%%%%%%%%%%%%%%%%%%%%%%%%%%%%%%%%%%%%%%%%%%%%%%%%%%%%%%%%%%
\subsection{Homotopy of Paths}
\label{subsec:homotopy-of-paths}
%%%%%%%%%%%%%%%%%%%%%%%%%%%%%%%%%%%%%%%%%%%%%%%%%%%%%%%%%%%%

Two paths

\[
\alpha,\beta:[0,1]\to X
\]

with the same endpoints may be homotopic while keeping those endpoints fixed.

This means there exists

\[
H:[0,1]\times[0,1]\to X
\]

such that

\[
H(s,0)=\alpha(s),
\]

\[
H(s,1)=\beta(s),
\]

and

\[
\boxed{
H(0,t)=\alpha(0)=\beta(0),
}
\]

\[
\boxed{
H(1,t)=\alpha(1)=\beta(1).
}
\]

This is called a \emph{homotopy relative to the endpoints}.

%%%%%%%%%%%%%%%%%%%%%%%%%%%%%%%%%%%%%%%%%%%%%%%%%%%%%%%%%%%%
\subsection{Homotopy Classes of Paths}
\label{subsec:homotopy-classes-paths}
%%%%%%%%%%%%%%%%%%%%%%%%%%%%%%%%%%%%%%%%%%%%%%%%%%%%%%%%%%%%

The homotopy class of a path \(\gamma\) relative to its endpoints is written

\[
\boxed{
[\gamma].
}
\]

This notation means that paths continuously deformable into \(\gamma\), with
endpoints fixed, are regarded as equivalent.

%%%%%%%%%%%%%%%%%%%%%%%%%%%%%%%%%%%%%%%%%%%%%%%%%%%%%%%%%%%%
\subsection{Homotopy Versus Homeomorphism}
\label{subsec:homotopy-vs-homeomorphism}
%%%%%%%%%%%%%%%%%%%%%%%%%%%%%%%%%%%%%%%%%%%%%%%%%%%%%%%%%%%%

Homeomorphism and homotopy are different notions.

A homeomorphism

\[
X\cong Y
\]

requires a continuous bijection with continuous inverse.

Homotopy equivalence is weaker.

Two spaces may fail to be homeomorphic but still have the same homotopy
type.

\begin{importantbox}
Homeomorphism means the spaces have exactly the same topological structure.

Homotopy equivalence means they have the same structure from the coarser
viewpoint of continuous deformation.
\end{importantbox}

%%%%%%%%%%%%%%%%%%%%%%%%%%%%%%%%%%%%%%%%%%%%%%%%%%%%%%%%%%%%
\subsection{Homotopy Equivalence}
\label{subsec:homotopy-equivalence}
%%%%%%%%%%%%%%%%%%%%%%%%%%%%%%%%%%%%%%%%%%%%%%%%%%%%%%%%%%%%

\begin{definitionbox}{Homotopy Equivalence}
Spaces \(X\) and \(Y\) are \emph{homotopy equivalent} if there exist
continuous maps

\[
f:X\to Y
\]

and

\[
g:Y\to X
\]

such that

\[
\boxed{
g\circ f\simeq\operatorname{id}_X
}
\]

and

\[
\boxed{
f\circ g\simeq\operatorname{id}_Y.
}
\]
\end{definitionbox}

We may write

\[
\boxed{
X\simeq Y.
}
\]

%%%%%%%%%%%%%%%%%%%%%%%%%%%%%%%%%%%%%%%%%%%%%%%%%%%%%%%%%%%%
\subsection{Contractible Spaces}
\label{subsec:contractible-spaces}
%%%%%%%%%%%%%%%%%%%%%%%%%%%%%%%%%%%%%%%%%%%%%%%%%%%%%%%%%%%%

\begin{definitionbox}{Contractible Space}
A space \(X\) is \emph{contractible} if it is homotopy equivalent to a
one-point space.
\end{definitionbox}

Equivalently, the identity map

\[
\operatorname{id}_X
\]

is homotopic to a constant map.

Intuitively, the entire space can be continuously contracted to one point.

%%%%%%%%%%%%%%%%%%%%%%%%%%%%%%%%%%%%%%%%%%%%%%%%%%%%%%%%%%%%
\subsection{Convex Sets Are Contractible}
\label{subsec:convex-contractible}
%%%%%%%%%%%%%%%%%%%%%%%%%%%%%%%%%%%%%%%%%%%%%%%%%%%%%%%%%%%%

Let

\[
C\subseteq\R^n
\]

be convex and choose

\[
x_0\in C.
\]

Define

\[
\boxed{
H(x,t)
=
(1-t)x+t x_0.
}
\]

Since \(C\) is convex,

\[
H(x,t)\in C.
\]

At \(t=0\),

\[
H(x,0)=x,
\]

while at \(t=1\),

\[
H(x,1)=x_0.
\]

Therefore

\[
\boxed{
C\text{ is contractible}.
}
\]

In particular,

\[
\R^n
\]

is contractible.

%%%%%%%%%%%%%%%%%%%%%%%%%%%%%%%%%%%%%%%%%%%%%%%%%%%%%%%%%%%%
\subsection{Retracts}
\label{subsec:retracts}
%%%%%%%%%%%%%%%%%%%%%%%%%%%%%%%%%%%%%%%%%%%%%%%%%%%%%%%%%%%%

\begin{definitionbox}{Retract}
Let

\[
A\subseteq X.
\]

A continuous map

\[
r:X\to A
\]

is a \emph{retraction} if

\[
\boxed{
r(a)=a
}
\]

for every

\[
a\in A.
\]
\end{definitionbox}

Then \(A\) is called a \emph{retract} of \(X\).

%%%%%%%%%%%%%%%%%%%%%%%%%%%%%%%%%%%%%%%%%%%%%%%%%%%%%%%%%%%%
\subsection{Deformation Retracts}
\label{subsec:deformation-retracts}
%%%%%%%%%%%%%%%%%%%%%%%%%%%%%%%%%%%%%%%%%%%%%%%%%%%%%%%%%%%%

\begin{definitionbox}{Deformation Retract}
A subspace \(A\subseteq X\) is a \emph{deformation retract} of \(X\) if
there exists a homotopy

\[
H:X\times[0,1]\to X
\]

such that

\[
H(x,0)=x,
\]

\[
H(x,1)\in A,
\]

and

\[
\boxed{
H(a,t)=a
}
\]

for all \(a\in A\) and all \(t\in[0,1]\).
\end{definitionbox}

A deformation retract has the same homotopy type as the original space.

%%%%%%%%%%%%%%%%%%%%%%%%%%%%%%%%%%%%%%%%%%%%%%%%%%%%%%%%%%%%
\subsection{Example: Punctured Plane}
\label{subsec:punctured-plane-deformation-retract}
%%%%%%%%%%%%%%%%%%%%%%%%%%%%%%%%%%%%%%%%%%%%%%%%%%%%%%%%%%%%

Consider

\[
\R^2\setminus\{0\}.
\]

Every nonzero point can be radially moved onto the unit circle.

Define

\[
r(x)=\frac{x}{\|x\|}.
\]

A deformation may be written

\[
\boxed{
H(x,t)
=
\left(
(1-t)+\frac{t}{\|x\|}
\right)x.
}
\]

Therefore

\[
\boxed{
\R^2\setminus\{0\}
\simeq
S^1.
}
\]

This example is extremely important because it reduces a two-dimensional
space to a much simpler one-dimensional space without changing its homotopy
type.

%%%%%%%%%%%%%%%%%%%%%%%%%%%%%%%%%%%%%%%%%%%%%%%%%%%%%%%%%%%%
\subsection{The Fundamental Group}
\label{subsec:fundamental-group}
%%%%%%%%%%%%%%%%%%%%%%%%%%%%%%%%%%%%%%%%%%%%%%%%%%%%%%%%%%%%

Let

\[
(X,x_0)
\]

be a topological space with chosen basepoint.

Consider loops based at \(x_0\).

Two loops are considered equivalent when they are homotopic relative to the
basepoint.

\begin{definitionbox}{Fundamental Group}
The \emph{fundamental group} of \(X\) at \(x_0\) is

\[
\boxed{
\pi_1(X,x_0)
=
\{
[\gamma]:
\gamma
\text{ is a loop based at }x_0
\}.
}
\]

The group operation is induced by path concatenation.
\end{definitionbox}

%%%%%%%%%%%%%%%%%%%%%%%%%%%%%%%%%%%%%%%%%%%%%%%%%%%%%%%%%%%%
\subsection{Group Structure of the Fundamental Group}
\label{subsec:fundamental-group-structure}
%%%%%%%%%%%%%%%%%%%%%%%%%%%%%%%%%%%%%%%%%%%%%%%%%%%%%%%%%%%%

The product is

\[
\boxed{
[\alpha][\beta]
=
[\alpha*\beta].
}
\]

The identity element is the class of the constant loop:

\[
\boxed{
[e_{x_0}].
}
\]

The inverse of

\[
[\gamma]
\]

is

\[
\boxed{
[\overline{\gamma}].
}
\]

Associativity holds on homotopy classes.

Thus

\[
\pi_1(X,x_0)
\]

is genuinely a group.

%%%%%%%%%%%%%%%%%%%%%%%%%%%%%%%%%%%%%%%%%%%%%%%%%%%%%%%%%%%%
\subsection{What the Fundamental Group Measures}
\label{subsec:fundamental-group-meaning}
%%%%%%%%%%%%%%%%%%%%%%%%%%%%%%%%%%%%%%%%%%%%%%%%%%%%%%%%%%%%

The fundamental group detects one-dimensional holes.

Very roughly,

\[
\boxed{
\pi_1(X)
\text{ records how loops can wind through }X.
}
\]

A loop that can be continuously contracted to a point represents the
identity element.

A loop that cannot be contracted may represent a nontrivial element.

%%%%%%%%%%%%%%%%%%%%%%%%%%%%%%%%%%%%%%%%%%%%%%%%%%%%%%%%%%%%
\subsection{Fundamental Group of a Point}
\label{subsec:fundamental-group-point}
%%%%%%%%%%%%%%%%%%%%%%%%%%%%%%%%%%%%%%%%%%%%%%%%%%%%%%%%%%%%

A one-point space has only the constant loop.

Therefore

\[
\boxed{
\pi_1(\{*\})
=
\{e\}.
}
\]

This is the trivial group.

%%%%%%%%%%%%%%%%%%%%%%%%%%%%%%%%%%%%%%%%%%%%%%%%%%%%%%%%%%%%
\subsection{Fundamental Group of a Contractible Space}
\label{subsec:fundamental-group-contractible}
%%%%%%%%%%%%%%%%%%%%%%%%%%%%%%%%%%%%%%%%%%%%%%%%%%%%%%%%%%%%

If \(X\) is contractible, then every loop contracts to a point.

Therefore

\[
\boxed{
\pi_1(X,x_0)
\cong
\{e\}.
}
\]

In particular,

\[
\boxed{
\pi_1(\R^n)
=
\{e\}.
}
\]

%%%%%%%%%%%%%%%%%%%%%%%%%%%%%%%%%%%%%%%%%%%%%%%%%%%%%%%%%%%%
\subsection{The Fundamental Group of the Circle}
\label{subsec:fundamental-group-circle}
%%%%%%%%%%%%%%%%%%%%%%%%%%%%%%%%%%%%%%%%%%%%%%%%%%%%%%%%%%%%

One of the central computations of algebraic topology is

\[
\boxed{
\pi_1(S^1)
\cong
\Z.
}
\]

The integer records the number of times a loop winds around the circle,
including orientation.

For example,

\[
0
\]

represents a contractible loop,

\[
1
\]

represents one counterclockwise winding,

\[
-1
\]

represents one clockwise winding, and

\[
n
\]

represents \(n\) signed windings.

%%%%%%%%%%%%%%%%%%%%%%%%%%%%%%%%%%%%%%%%%%%%%%%%%%%%%%%%%%%%
\subsection{Winding Number}
\label{subsec:topological-winding-number}
%%%%%%%%%%%%%%%%%%%%%%%%%%%%%%%%%%%%%%%%%%%%%%%%%%%%%%%%%%%%

A standard loop on the circle is

\[
\boxed{
\gamma_n(t)
=
e^{2\pi i n t}.
}
\]

Here

\[
i^2=-1.
\]

As \(t\) runs from \(0\) to \(1\), the loop winds around the circle \(n\)
times.

Its homotopy class corresponds to

\[
n\in\Z.
\]

Thus the group operation corresponds to addition:

\[
\boxed{
[\gamma_m][\gamma_n]
=
[\gamma_{m+n}].
}
\]

%%%%%%%%%%%%%%%%%%%%%%%%%%%%%%%%%%%%%%%%%%%%%%%%%%%%%%%%%%%%
\subsection{Fundamental Group of the Punctured Plane}
\label{subsec:fundamental-group-punctured-plane}
%%%%%%%%%%%%%%%%%%%%%%%%%%%%%%%%%%%%%%%%%%%%%%%%%%%%%%%%%%%%

Since

\[
\R^2\setminus\{0\}
\simeq
S^1,
\]

homotopy invariance gives

\[
\boxed{
\pi_1(\R^2\setminus\{0\})
\cong
\Z.
}
\]

The missing point creates a topological hole around which loops may wind.

%%%%%%%%%%%%%%%%%%%%%%%%%%%%%%%%%%%%%%%%%%%%%%%%%%%%%%%%%%%%
\subsection{Simply Connected Spaces}
\label{subsec:simply-connected-spaces}
%%%%%%%%%%%%%%%%%%%%%%%%%%%%%%%%%%%%%%%%%%%%%%%%%%%%%%%%%%%%

\begin{definitionbox}{Simply Connected Space}
A topological space \(X\) is \emph{simply connected} if:

\begin{enumerate}
    \item \(X\) is path connected; and
    \item
    \[
    \boxed{
    \pi_1(X,x_0)
    }
    \]
    is trivial.
\end{enumerate}
\end{definitionbox}

Equivalently, every loop can be continuously contracted to a point.

%%%%%%%%%%%%%%%%%%%%%%%%%%%%%%%%%%%%%%%%%%%%%%%%%%%%%%%%%%%%
\subsection{Examples of Simply Connected Spaces}
\label{subsec:simply-connected-examples}
%%%%%%%%%%%%%%%%%%%%%%%%%%%%%%%%%%%%%%%%%%%%%%%%%%%%%%%%%%%%

Examples include

\[
\boxed{
\R^n,
}
\]

\[
\boxed{
D^n,
}
\]

and

\[
\boxed{
S^n
\quad\text{for }n\geq2.
}
\]

The circle is not simply connected because

\[
\pi_1(S^1)\cong\Z.
\]

%%%%%%%%%%%%%%%%%%%%%%%%%%%%%%%%%%%%%%%%%%%%%%%%%%%%%%%%%%%%
\subsection{Basepoints and Path-Connected Spaces}
\label{subsec:basepoint-fundamental-group}
%%%%%%%%%%%%%%%%%%%%%%%%%%%%%%%%%%%%%%%%%%%%%%%%%%%%%%%%%%%%

The fundamental group formally depends on a basepoint:

\[
\pi_1(X,x_0).
\]

However, if \(X\) is path connected and

\[
x_0,x_1\in X,
\]

then

\[
\boxed{
\pi_1(X,x_0)
\cong
\pi_1(X,x_1).
}
\]

The isomorphism depends on a choice of path between the basepoints.

Therefore one often writes simply

\[
\pi_1(X)
\]

when only the group up to isomorphism matters.

%%%%%%%%%%%%%%%%%%%%%%%%%%%%%%%%%%%%%%%%%%%%%%%%%%%%%%%%%%%%
\subsection{Induced Homomorphisms}
\label{subsec:induced-homomorphisms-pi1}
%%%%%%%%%%%%%%%%%%%%%%%%%%%%%%%%%%%%%%%%%%%%%%%%%%%%%%%%%%%%

Let

\[
f:(X,x_0)\to(Y,y_0)
\]

be continuous with

\[
f(x_0)=y_0.
\]

A loop

\[
\gamma:[0,1]\to X
\]

produces a loop

\[
f\circ\gamma
\]

in \(Y\).

This defines a homomorphism

\[
\boxed{
f_*:
\pi_1(X,x_0)
\to
\pi_1(Y,y_0)
}
\]

by

\[
\boxed{
f_*([\gamma])
=
[f\circ\gamma].
}
\]

The symbol

\[
f_*
\]

is read ``\(f\) star.''

%%%%%%%%%%%%%%%%%%%%%%%%%%%%%%%%%%%%%%%%%%%%%%%%%%%%%%%%%%%%
\subsection{Functorial Behavior}
\label{subsec:fundamental-group-functoriality}
%%%%%%%%%%%%%%%%%%%%%%%%%%%%%%%%%%%%%%%%%%%%%%%%%%%%%%%%%%%%

The induced homomorphisms satisfy

\[
\boxed{
(\operatorname{id}_X)_*
=
\operatorname{id}_{\pi_1(X)}
}
\]

and

\[
\boxed{
(g\circ f)_*
=
g_*\circ f_*.
}
\]

This is an example of \emph{functoriality}.

Topology provides spaces and continuous maps.

Algebra provides groups and homomorphisms.

The fundamental-group construction respects composition.

%%%%%%%%%%%%%%%%%%%%%%%%%%%%%%%%%%%%%%%%%%%%%%%%%%%%%%%%%%%%
\subsection{Homotopy Invariance}
\label{subsec:fundamental-group-homotopy-invariance}
%%%%%%%%%%%%%%%%%%%%%%%%%%%%%%%%%%%%%%%%%%%%%%%%%%%%%%%%%%%%

Homotopy-equivalent spaces have isomorphic fundamental groups.

Thus

\[
\boxed{
X\simeq Y
\Longrightarrow
\pi_1(X)\cong\pi_1(Y)
}
\]

under the usual path-connected and basepoint considerations.

Therefore the fundamental group is a homotopy invariant.

%%%%%%%%%%%%%%%%%%%%%%%%%%%%%%%%%%%%%%%%%%%%%%%%%%%%%%%%%%%%
\subsection{The Fundamental Group as an Obstruction}
\label{subsec:fundamental-group-obstruction}
%%%%%%%%%%%%%%%%%%%%%%%%%%%%%%%%%%%%%%%%%%%%%%%%%%%%%%%%%%%%

Suppose

\[
\pi_1(X)\not\cong\pi_1(Y).
\]

Then \(X\) and \(Y\) cannot be homotopy equivalent.

Consequently, they cannot be homeomorphic.

For example,

\[
\pi_1(\R^2)=\{e\}
\]

while

\[
\pi_1(\R^2\setminus\{0\})\cong\Z.
\]

Hence

\[
\boxed{
\R^2
\not\simeq
\R^2\setminus\{0\}.
}
\]

%%%%%%%%%%%%%%%%%%%%%%%%%%%%%%%%%%%%%%%%%%%%%%%%%%%%%%%%%%%%
\subsection{The Torus}
\label{subsec:fundamental-group-torus}
%%%%%%%%%%%%%%%%%%%%%%%%%%%%%%%%%%%%%%%%%%%%%%%%%%%%%%%%%%%%

The torus is

\[
\boxed{
T^2=S^1\times S^1.
}
\]

Its fundamental group is

\[
\boxed{
\pi_1(T^2)
\cong
\Z\times\Z.
}
\]

One generator corresponds to traveling around one circular direction.

The other corresponds to traveling around the second circular direction.

The two generators commute.

%%%%%%%%%%%%%%%%%%%%%%%%%%%%%%%%%%%%%%%%%%%%%%%%%%%%%%%%%%%%
\subsection{Wedge of Circles}
\label{subsec:wedge-of-circles}
%%%%%%%%%%%%%%%%%%%%%%%%%%%%%%%%%%%%%%%%%%%%%%%%%%%%%%%%%%%%

The wedge

\[
S^1\vee S^1
\]

is obtained by identifying one point of one circle with one point of another.

The symbol

\[
\vee
\]

is read ``wedge.''

Its fundamental group is

\[
\boxed{
\pi_1(S^1\vee S^1)
\cong
F_2,
}
\]

where \(F_2\) is the free group on two generators.

Unlike

\[
\Z\times\Z,
\]

the free group \(F_2\) is nonabelian.

%%%%%%%%%%%%%%%%%%%%%%%%%%%%%%%%%%%%%%%%%%%%%%%%%%%%%%%%%%%%
\subsection{The Seifert--van Kampen Principle}
\label{subsec:seifert-van-kampen}
%%%%%%%%%%%%%%%%%%%%%%%%%%%%%%%%%%%%%%%%%%%%%%%%%%%%%%%%%%%%

The fundamental group of a complicated space can often be computed by
decomposing the space into simpler overlapping pieces.

The main theorem governing this process is the
\emph{Seifert--van Kampen theorem}.

In simplified form, if

\[
X=U\cup V
\]

with suitable path-connected open sets \(U,V\), then

\[
\pi_1(X)
\]

can be assembled from

\[
\pi_1(U),
\qquad
\pi_1(V),
\qquad
\pi_1(U\cap V).
\]

Algebraically, this construction is expressed using a pushout, often
realized as an amalgamated free product.

%%%%%%%%%%%%%%%%%%%%%%%%%%%%%%%%%%%%%%%%%%%%%%%%%%%%%%%%%%%%
\subsection{Fundamental Groups of Surfaces}
\label{subsec:fundamental-groups-surfaces}
%%%%%%%%%%%%%%%%%%%%%%%%%%%%%%%%%%%%%%%%%%%%%%%%%%%%%%%%%%%%

Fundamental groups distinguish many surfaces.

For example,

\[
\boxed{
\pi_1(S^2)=\{e\},
}
\]

while

\[
\boxed{
\pi_1(T^2)
\cong
\Z^2.
}
\]

For an orientable closed surface of genus \(g\),

\[
\boxed{
\pi_1(\Sigma_g)
=
\left\langle
a_1,b_1,\ldots,a_g,b_g
\;\middle|\;
[a_1,b_1]\cdots[a_g,b_g]=e
\right\rangle.
}
\]

Here

\[
[a,b]
=
aba^{-1}b^{-1}
\]

is the group commutator.

%%%%%%%%%%%%%%%%%%%%%%%%%%%%%%%%%%%%%%%%%%%%%%%%%%%%%%%%%%%%
\subsection{Covering Spaces}
\label{subsec:covering-spaces}
%%%%%%%%%%%%%%%%%%%%%%%%%%%%%%%%%%%%%%%%%%%%%%%%%%%%%%%%%%%%

\begin{definitionbox}{Covering Map}
A continuous surjective map

\[
p:\widetilde{X}\to X
\]

is a \emph{covering map} if every point \(x\in X\) has an open neighborhood
\(U\) such that

\[
p^{-1}(U)
\]

is a disjoint union of open sets, each of which is mapped homeomorphically
onto \(U\) by \(p\).
\end{definitionbox}

The space

\[
\widetilde{X}
\]

is called a \emph{covering space} of \(X\).

The symbol

\[
\widetilde{X}
\]

is read ``X tilde.''

%%%%%%%%%%%%%%%%%%%%%%%%%%%%%%%%%%%%%%%%%%%%%%%%%%%%%%%%%%%%
\subsection{The Real Line Covers the Circle}
\label{subsec:real-line-covers-circle}
%%%%%%%%%%%%%%%%%%%%%%%%%%%%%%%%%%%%%%%%%%%%%%%%%%%%%%%%%%%%

Consider

\[
p:\R\to S^1
\]

defined by

\[
\boxed{
p(t)=e^{2\pi i t}.
}
\]

Then

\[
p(t+n)=p(t)
\]

for every

\[
n\in\Z.
\]

Locally, \(p\) looks like a homeomorphism.

Globally, infinitely many points of \(\R\) map to the same point of \(S^1\).

Thus

\[
\boxed{
\R\to S^1
}
\]

is a covering map.

%%%%%%%%%%%%%%%%%%%%%%%%%%%%%%%%%%%%%%%%%%%%%%%%%%%%%%%%%%%%
\subsection{Fibers of Covering Maps}
\label{subsec:covering-fibers}
%%%%%%%%%%%%%%%%%%%%%%%%%%%%%%%%%%%%%%%%%%%%%%%%%%%%%%%%%%%%

For

\[
x\in X,
\]

the set

\[
\boxed{
p^{-1}(x)
}
\]

is called the \emph{fiber} over \(x\).

For the covering

\[
p(t)=e^{2\pi i t},
\]

the fiber over

\[
1\in S^1
\]

is

\[
\boxed{
p^{-1}(1)=\Z.
}
\]

%%%%%%%%%%%%%%%%%%%%%%%%%%%%%%%%%%%%%%%%%%%%%%%%%%%%%%%%%%%%
\subsection{Path Lifting}
\label{subsec:path-lifting}
%%%%%%%%%%%%%%%%%%%%%%%%%%%%%%%%%%%%%%%%%%%%%%%%%%%%%%%%%%%%

One of the key properties of covering maps is that paths can be lifted.

Suppose

\[
p:\widetilde{X}\to X
\]

is a covering map and

\[
\gamma:[0,1]\to X
\]

is a path.

Choose

\[
\widetilde{x}_0\in p^{-1}(\gamma(0)).
\]

Then, under the covering-space hypotheses, there exists a unique path

\[
\widetilde{\gamma}:[0,1]\to\widetilde{X}
\]

such that

\[
\boxed{
p\circ\widetilde{\gamma}
=
\gamma
}
\]

and

\[
\widetilde{\gamma}(0)=\widetilde{x}_0.
\]

%%%%%%%%%%%%%%%%%%%%%%%%%%%%%%%%%%%%%%%%%%%%%%%%%%%%%%%%%%%%
\subsection{Homotopy Lifting}
\label{subsec:homotopy-lifting}
%%%%%%%%%%%%%%%%%%%%%%%%%%%%%%%%%%%%%%%%%%%%%%%%%%%%%%%%%%%%

Covering maps also satisfy a homotopy lifting property.

Under appropriate hypotheses, a homotopy in the base space can be lifted
once an initial lift is specified.

This allows topological questions about \(X\) to be studied in the often
simpler space

\[
\widetilde{X}.
\]

%%%%%%%%%%%%%%%%%%%%%%%%%%%%%%%%%%%%%%%%%%%%%%%%%%%%%%%%%%%%
\subsection{Universal Covers}
\label{subsec:universal-covers}
%%%%%%%%%%%%%%%%%%%%%%%%%%%%%%%%%%%%%%%%%%%%%%%%%%%%%%%%%%%%

\begin{definitionbox}{Universal Cover}
A \emph{universal covering space} of \(X\) is a simply connected covering
space

\[
\boxed{
p:\widetilde{X}\to X.
}
\]
\end{definitionbox}

For example,

\[
\boxed{
\R
}
\]

is the universal cover of

\[
\boxed{
S^1.
}
\]

Under standard local hypotheses, connected spaces admit universal covers.

%%%%%%%%%%%%%%%%%%%%%%%%%%%%%%%%%%%%%%%%%%%%%%%%%%%%%%%%%%%%
\subsection{Deck Transformations}
\label{subsec:deck-transformations}
%%%%%%%%%%%%%%%%%%%%%%%%%%%%%%%%%%%%%%%%%%%%%%%%%%%%%%%%%%%%

A \emph{deck transformation} of a covering

\[
p:\widetilde{X}\to X
\]

is a homeomorphism

\[
h:\widetilde{X}\to\widetilde{X}
\]

satisfying

\[
\boxed{
p\circ h=p.
}
\]

For the universal covering

\[
\R\to S^1,
\]

the transformations

\[
\boxed{
h_n(t)=t+n,
\qquad
n\in\Z,
}
\]

are deck transformations.

Their group is isomorphic to

\[
\Z.
\]

%%%%%%%%%%%%%%%%%%%%%%%%%%%%%%%%%%%%%%%%%%%%%%%%%%%%%%%%%%%%
\subsection{Coverings and the Fundamental Group}
\label{subsec:coverings-fundamental-group}
%%%%%%%%%%%%%%%%%%%%%%%%%%%%%%%%%%%%%%%%%%%%%%%%%%%%%%%%%%%%

Under suitable hypotheses, connected covering spaces of a space \(X\)
correspond to subgroups of

\[
\pi_1(X).
\]

This creates a deep relationship:

\[
\boxed{
\text{geometry of covering spaces}
\longleftrightarrow
\text{algebra of subgroups}.
}
\]

For universal covers, the deck-transformation group is closely related to
the fundamental group.

%%%%%%%%%%%%%%%%%%%%%%%%%%%%%%%%%%%%%%%%%%%%%%%%%%%%%%%%%%%%
\subsection{Higher Homotopy Groups}
\label{subsec:higher-homotopy-groups}
%%%%%%%%%%%%%%%%%%%%%%%%%%%%%%%%%%%%%%%%%%%%%%%%%%%%%%%%%%%%

The fundamental group studies maps from the circle

\[
S^1.
\]

This idea extends to higher-dimensional spheres.

The \(n\)-th homotopy group is denoted

\[
\boxed{
\pi_n(X,x_0).
}
\]

Roughly, it studies based maps

\[
S^n\to X
\]

up to based homotopy.

For

\[
n\geq2,
\]

the groups

\[
\pi_n(X)
\]

are abelian.

%%%%%%%%%%%%%%%%%%%%%%%%%%%%%%%%%%%%%%%%%%%%%%%%%%%%%%%%%%%%
\subsection{Homotopy Groups of Spheres}
\label{subsec:homotopy-groups-spheres}
%%%%%%%%%%%%%%%%%%%%%%%%%%%%%%%%%%%%%%%%%%%%%%%%%%%%%%%%%%%%

Some basic examples are

\[
\boxed{
\pi_n(S^n)\cong\Z
}
\]

for

\[
n\geq1.
\]

However, groups of the form

\[
\pi_k(S^n)
\]

for general \(k\) and \(n\) can be extremely difficult to compute.

This difficulty is one reason homology theory is so useful.

%%%%%%%%%%%%%%%%%%%%%%%%%%%%%%%%%%%%%%%%%%%%%%%%%%%%%%%%%%%%
\subsection{Why Homology?}
\label{subsec:why-homology}
%%%%%%%%%%%%%%%%%%%%%%%%%%%%%%%%%%%%%%%%%%%%%%%%%%%%%%%%%%%%

The fundamental group is powerful, but it is not always easy to compute.

Homology provides a sequence of abelian groups

\[
\boxed{
H_0(X),
H_1(X),
H_2(X),
\ldots
}
\]

that detect holes in different dimensions.

Very roughly:

\[
\boxed{
H_0
\longleftrightarrow
\text{connected components},
}
\]

\[
\boxed{
H_1
\longleftrightarrow
\text{one-dimensional holes},
}
\]

\[
\boxed{
H_2
\longleftrightarrow
\text{two-dimensional cavities},
}
\]

and so forth.

This interpretation is useful intuition, though the formal definition is
algebraic.

%%%%%%%%%%%%%%%%%%%%%%%%%%%%%%%%%%%%%%%%%%%%%%%%%%%%%%%%%%%%
\subsection{Simplices}
\label{subsec:simplices}
%%%%%%%%%%%%%%%%%%%%%%%%%%%%%%%%%%%%%%%%%%%%%%%%%%%%%%%%%%%%

A simplex is the higher-dimensional analogue of a triangle.

A \(0\)-simplex is a point.

A \(1\)-simplex is a line segment.

A \(2\)-simplex is a triangle.

A \(3\)-simplex is a tetrahedron.

In general, an \(n\)-simplex is the convex hull of \(n+1\) affinely
independent points.

%%%%%%%%%%%%%%%%%%%%%%%%%%%%%%%%%%%%%%%%%%%%%%%%%%%%%%%%%%%%
\subsection{Standard Simplex}
\label{subsec:standard-simplex}
%%%%%%%%%%%%%%%%%%%%%%%%%%%%%%%%%%%%%%%%%%%%%%%%%%%%%%%%%%%%

The standard \(n\)-simplex is

\[
\boxed{
\Delta^n
=
\left\{
(t_0,\ldots,t_n)\in\R^{n+1}:
t_i\geq0,
\;
\sum_{i=0}^{n}t_i=1
\right\}.
}
\]

The uppercase Greek letter

\[
\Delta
\]

is read ``delta.''

%%%%%%%%%%%%%%%%%%%%%%%%%%%%%%%%%%%%%%%%%%%%%%%%%%%%%%%%%%%%
\subsection{Faces of a Simplex}
\label{subsec:faces-simplex}
%%%%%%%%%%%%%%%%%%%%%%%%%%%%%%%%%%%%%%%%%%%%%%%%%%%%%%%%%%%%

A face of a simplex is obtained by omitting one or more vertices.

For an oriented \(2\)-simplex

\[
[v_0,v_1,v_2],
\]

its oriented edges are

\[
[v_1,v_2],
\qquad
[v_0,v_2],
\qquad
[v_0,v_1].
\]

Signs will be used to encode their orientations.

%%%%%%%%%%%%%%%%%%%%%%%%%%%%%%%%%%%%%%%%%%%%%%%%%%%%%%%%%%%%
\subsection{Simplicial Complexes}
\label{subsec:simplicial-complexes}
%%%%%%%%%%%%%%%%%%%%%%%%%%%%%%%%%%%%%%%%%%%%%%%%%%%%%%%%%%%%

\begin{definitionbox}{Simplicial Complex}
A \emph{simplicial complex} is a collection of simplices satisfying:

\begin{enumerate}
    \item every face of a simplex in the collection is also in the
    collection;

    \item the intersection of any two simplices is either empty or a common
    face of both.
\end{enumerate}
\end{definitionbox}

Simplicial complexes allow spaces to be represented combinatorially.

%%%%%%%%%%%%%%%%%%%%%%%%%%%%%%%%%%%%%%%%%%%%%%%%%%%%%%%%%%%%
\subsection{Triangulation}
\label{subsec:topological-triangulation}
%%%%%%%%%%%%%%%%%%%%%%%%%%%%%%%%%%%%%%%%%%%%%%%%%%%%%%%%%%%%

A triangulation represents a topological space using a simplicial complex.

For example, a circle can be represented by a polygonal loop.

A sphere can be represented by the boundary of a tetrahedron.

This turns geometric structure into finite combinatorial data.

%%%%%%%%%%%%%%%%%%%%%%%%%%%%%%%%%%%%%%%%%%%%%%%%%%%%%%%%%%%%
\subsection{Oriented Simplices}
\label{subsec:oriented-simplices}
%%%%%%%%%%%%%%%%%%%%%%%%%%%%%%%%%%%%%%%%%%%%%%%%%%%%%%%%%%%%

An orientation of a simplex is determined by an ordering of its vertices,
up to even permutation.

For example,

\[
[v_0,v_1]
\]

and

\[
[v_1,v_0]
\]

have opposite orientations:

\[
\boxed{
[v_1,v_0]
=
-[v_0,v_1].
}
\]

Orientation signs are essential in defining boundary operators.

%%%%%%%%%%%%%%%%%%%%%%%%%%%%%%%%%%%%%%%%%%%%%%%%%%%%%%%%%%%%
\subsection{Chain Groups}
\label{subsec:chain-groups}
%%%%%%%%%%%%%%%%%%%%%%%%%%%%%%%%%%%%%%%%%%%%%%%%%%%%%%%%%%%%

Let \(K\) be a simplicial complex.

The \(n\)-th chain group

\[
\boxed{
C_n(K)
}
\]

is the free abelian group generated by the oriented \(n\)-simplices of
\(K\).

Thus an \(n\)-chain has the form

\[
\boxed{
c
=
\sum_i a_i\sigma_i,
\qquad
a_i\in\Z,
}
\]

where the \(\sigma_i\) are oriented \(n\)-simplices.

The lowercase Greek letter

\[
\sigma
\]

is read ``sigma.''

%%%%%%%%%%%%%%%%%%%%%%%%%%%%%%%%%%%%%%%%%%%%%%%%%%%%%%%%%%%%
\subsection{Boundary Operator}
\label{subsec:boundary-operator}
%%%%%%%%%%%%%%%%%%%%%%%%%%%%%%%%%%%%%%%%%%%%%%%%%%%%%%%%%%%%

The boundary homomorphism is

\[
\boxed{
\partial_n:
C_n(K)\to C_{n-1}(K).
}
\]

For an oriented simplex

\[
[v_0,\ldots,v_n],
\]

define

\[
\boxed{
\partial_n[v_0,\ldots,v_n]
=
\sum_{i=0}^{n}
(-1)^i
[v_0,\ldots,\widehat{v_i},\ldots,v_n].
}
\]

The notation

\[
\widehat{v_i}
\]

means that \(v_i\) is omitted.

%%%%%%%%%%%%%%%%%%%%%%%%%%%%%%%%%%%%%%%%%%%%%%%%%%%%%%%%%%%%
\subsection{Boundary of an Edge}
\label{subsec:boundary-edge}
%%%%%%%%%%%%%%%%%%%%%%%%%%%%%%%%%%%%%%%%%%%%%%%%%%%%%%%%%%%%

For an oriented edge

\[
[v_0,v_1],
\]

the boundary is

\[
\boxed{
\partial_1[v_0,v_1]
=
[v_1]-[v_0].
}
\]

Thus the boundary consists of the terminal vertex minus the initial vertex.

%%%%%%%%%%%%%%%%%%%%%%%%%%%%%%%%%%%%%%%%%%%%%%%%%%%%%%%%%%%%
\subsection{Boundary of a Triangle}
\label{subsec:boundary-triangle}
%%%%%%%%%%%%%%%%%%%%%%%%%%%%%%%%%%%%%%%%%%%%%%%%%%%%%%%%%%%%

For an oriented triangle

\[
[v_0,v_1,v_2],
\]

we obtain

\[
\boxed{
\partial_2[v_0,v_1,v_2]
=
[v_1,v_2]
-
[v_0,v_2]
+
[v_0,v_1].
}
\]

The signs ensure that adjacent faces cancel correctly when simplices are
assembled.

%%%%%%%%%%%%%%%%%%%%%%%%%%%%%%%%%%%%%%%%%%%%%%%%%%%%%%%%%%%%
\subsection{Boundary of a Boundary}
\label{subsec:boundary-squared-zero}
%%%%%%%%%%%%%%%%%%%%%%%%%%%%%%%%%%%%%%%%%%%%%%%%%%%%%%%%%%%%

The fundamental algebraic property is

\[
\boxed{
\partial_{n-1}\circ\partial_n=0.
}
\]

This is often abbreviated as

\[
\boxed{
\partial^2=0.
}
\]

In words:

\[
\boxed{
\text{the boundary of a boundary is zero}.
}
\]

%%%%%%%%%%%%%%%%%%%%%%%%%%%%%%%%%%%%%%%%%%%%%%%%%%%%%%%%%%%%
\subsection{Worked Example: Boundary of a Triangle}
\label{subsec:worked-boundary-triangle}
%%%%%%%%%%%%%%%%%%%%%%%%%%%%%%%%%%%%%%%%%%%%%%%%%%%%%%%%%%%%

\begin{workedexamplebox}{Verify That the Boundary of a Boundary Is Zero}

Let

\[
\sigma=[v_0,v_1,v_2].
\]

Then

\[
\partial_2\sigma
=
[v_1,v_2]
-
[v_0,v_2]
+
[v_0,v_1].
\]

Apply \(\partial_1\):

\[
\partial_1\partial_2\sigma
=
([v_2]-[v_1])
-
([v_2]-[v_0])
+
([v_1]-[v_0]).
\]

Collect terms:

\[
[v_2]-[v_1]-[v_2]+[v_0]+[v_1]-[v_0]
=
0.
\]

\begin{answerbox}
\[
\boxed{
\partial_1\partial_2\sigma=0.
}
\]
\end{answerbox}
\end{workedexamplebox}

%%%%%%%%%%%%%%%%%%%%%%%%%%%%%%%%%%%%%%%%%%%%%%%%%%%%%%%%%%%%
\subsection{Chain Complexes}
\label{subsec:chain-complexes-topology}
%%%%%%%%%%%%%%%%%%%%%%%%%%%%%%%%%%%%%%%%%%%%%%%%%%%%%%%%%%%%

The chain groups and boundary maps form a sequence

\[
\boxed{
\cdots
\longrightarrow
C_{n+1}
\xrightarrow{\partial_{n+1}}
C_n
\xrightarrow{\partial_n}
C_{n-1}
\longrightarrow
\cdots
\longrightarrow
C_0
\longrightarrow
0.
}
\]

Because

\[
\partial_n\partial_{n+1}=0,
\]

this sequence is called a \emph{chain complex}.

%%%%%%%%%%%%%%%%%%%%%%%%%%%%%%%%%%%%%%%%%%%%%%%%%%%%%%%%%%%%
\subsection{Cycles}
\label{subsec:cycles}
%%%%%%%%%%%%%%%%%%%%%%%%%%%%%%%%%%%%%%%%%%%%%%%%%%%%%%%%%%%%

\begin{definitionbox}{Cycle}
An \(n\)-chain \(c\in C_n\) is an \emph{\(n\)-cycle} if

\[
\boxed{
\partial_n c=0.
}
\]
\end{definitionbox}

The group of \(n\)-cycles is

\[
\boxed{
Z_n
=
\ker(\partial_n).
}
\]

The letter \(Z\) is traditionally used for cycles.

%%%%%%%%%%%%%%%%%%%%%%%%%%%%%%%%%%%%%%%%%%%%%%%%%%%%%%%%%%%%
\subsection{Boundaries}
\label{subsec:homological-boundaries}
%%%%%%%%%%%%%%%%%%%%%%%%%%%%%%%%%%%%%%%%%%%%%%%%%%%%%%%%%%%%

\begin{definitionbox}{Boundary}
An \(n\)-chain \(c\in C_n\) is an \emph{\(n\)-boundary} if there exists an
\((n+1)\)-chain \(b\) such that

\[
\boxed{
c=\partial_{n+1}b.
}
\]
\end{definitionbox}

The group of \(n\)-boundaries is

\[
\boxed{
B_n
=
\operatorname{im}(\partial_{n+1}).
}
\]

Since

\[
\partial^2=0,
\]

every boundary is automatically a cycle:

\[
\boxed{
B_n\subseteq Z_n.
}
\]

%%%%%%%%%%%%%%%%%%%%%%%%%%%%%%%%%%%%%%%%%%%%%%%%%%%%%%%%%%%%
\subsection{Homology Groups}
\label{subsec:homology-groups}
%%%%%%%%%%%%%%%%%%%%%%%%%%%%%%%%%%%%%%%%%%%%%%%%%%%%%%%%%%%%

\begin{definitionbox}{Homology Group}
The \(n\)-th homology group is

\[
\boxed{
H_n
=
\frac{Z_n}{B_n}
=
\frac{
\ker(\partial_n)
}{
\operatorname{im}(\partial_{n+1})
}.
}
\]
\end{definitionbox}

Thus homology measures cycles that are not boundaries.

Conceptually,

\[
\boxed{
\text{homology}
=
\frac{
\text{closed cycles}
}{
\text{cycles bounding higher-dimensional pieces}
}.
}
\]

%%%%%%%%%%%%%%%%%%%%%%%%%%%%%%%%%%%%%%%%%%%%%%%%%%%%%%%%%%%%
\subsection{Geometric Meaning of Homology}
\label{subsec:geometric-meaning-homology}
%%%%%%%%%%%%%%%%%%%%%%%%%%%%%%%%%%%%%%%%%%%%%%%%%%%%%%%%%%%%

Suppose a loop forms the boundary of a filled disk.

Then that loop is a boundary and represents zero in first homology.

If a loop surrounds a hole and cannot be filled by a two-dimensional chain
inside the space, it may represent a nonzero homology class.

Thus homology detects failures of cycles to bound.

%%%%%%%%%%%%%%%%%%%%%%%%%%%%%%%%%%%%%%%%%%%%%%%%%%%%%%%%%%%%
\subsection{Zeroth Homology}
\label{subsec:zeroth-homology}
%%%%%%%%%%%%%%%%%%%%%%%%%%%%%%%%%%%%%%%%%%%%%%%%%%%%%%%%%%%%

The zeroth homology group records path components.

If \(X\) has \(k\) path components, then under standard hypotheses,

\[
\boxed{
H_0(X;\Z)
\cong
\Z^k.
}
\]

In particular, if \(X\) is nonempty and path connected,

\[
\boxed{
H_0(X;\Z)\cong\Z.
}
\]

%%%%%%%%%%%%%%%%%%%%%%%%%%%%%%%%%%%%%%%%%%%%%%%%%%%%%%%%%%%%
\subsection{Homology of a Point}
\label{subsec:homology-point}
%%%%%%%%%%%%%%%%%%%%%%%%%%%%%%%%%%%%%%%%%%%%%%%%%%%%%%%%%%%%

For a one-point space,

\[
\boxed{
H_0(\{*\};\Z)\cong\Z,
}
\]

and

\[
\boxed{
H_n(\{*\};\Z)=0
\qquad
(n\geq1).
}
\]

%%%%%%%%%%%%%%%%%%%%%%%%%%%%%%%%%%%%%%%%%%%%%%%%%%%%%%%%%%%%
\subsection{Homology of the Circle}
\label{subsec:homology-circle}
%%%%%%%%%%%%%%%%%%%%%%%%%%%%%%%%%%%%%%%%%%%%%%%%%%%%%%%%%%%%

For the circle,

\[
\boxed{
H_0(S^1;\Z)\cong\Z,
}
\]

\[
\boxed{
H_1(S^1;\Z)\cong\Z,
}
\]

and

\[
\boxed{
H_n(S^1;\Z)=0
\qquad
n\geq2.
}
\]

The nonzero group

\[
H_1(S^1)\cong\Z
\]

detects the one-dimensional hole in the circle.

%%%%%%%%%%%%%%%%%%%%%%%%%%%%%%%%%%%%%%%%%%%%%%%%%%%%%%%%%%%%
\subsection{Homology of the Sphere}
\label{subsec:homology-sphere}
%%%%%%%%%%%%%%%%%%%%%%%%%%%%%%%%%%%%%%%%%%%%%%%%%%%%%%%%%%%%

For the \(n\)-sphere \(S^n\),

\[
\boxed{
H_k(S^n;\Z)
\cong
\begin{cases}
\Z,
&
k=0,n,
\\[4pt]
0,
&
\text{otherwise}.
\end{cases}
}
\]

Thus \(S^n\) has a nontrivial \(n\)-dimensional homology class.

%%%%%%%%%%%%%%%%%%%%%%%%%%%%%%%%%%%%%%%%%%%%%%%%%%%%%%%%%%%%
\subsection{Homology of the Torus}
\label{subsec:homology-torus}
%%%%%%%%%%%%%%%%%%%%%%%%%%%%%%%%%%%%%%%%%%%%%%%%%%%%%%%%%%%%

For the torus \(T^2\),

\[
\boxed{
H_0(T^2;\Z)\cong\Z,
}
\]

\[
\boxed{
H_1(T^2;\Z)\cong\Z^2,
}
\]

\[
\boxed{
H_2(T^2;\Z)\cong\Z,
}
\]

and

\[
H_n(T^2;\Z)=0
\]

for \(n>2\).

The two generators of \(H_1\) correspond to the two independent circular
directions of the torus.

%%%%%%%%%%%%%%%%%%%%%%%%%%%%%%%%%%%%%%%%%%%%%%%%%%%%%%%%%%%%
\subsection{Homology of Contractible Spaces}
\label{subsec:homology-contractible}
%%%%%%%%%%%%%%%%%%%%%%%%%%%%%%%%%%%%%%%%%%%%%%%%%%%%%%%%%%%%

Homotopy-equivalent spaces have isomorphic homology groups.

Therefore a contractible nonempty space has the same homology as a point:

\[
\boxed{
H_0(X;\Z)\cong\Z,
}
\]

and

\[
\boxed{
H_n(X;\Z)=0
\qquad
n\geq1.
}
\]

%%%%%%%%%%%%%%%%%%%%%%%%%%%%%%%%%%%%%%%%%%%%%%%%%%%%%%%%%%%%
\subsection{Reduced Homology}
\label{subsec:reduced-homology}
%%%%%%%%%%%%%%%%%%%%%%%%%%%%%%%%%%%%%%%%%%%%%%%%%%%%%%%%%%%%

Reduced homology modifies the zeroth homology group so that a one-point
space has zero homology in every dimension.

It is denoted

\[
\boxed{
\widetilde{H}_n(X).
}
\]

The symbol

\[
\widetilde{\phantom{H}}
\]

is read ``tilde.''

For a nonempty path-connected space,

\[
\boxed{
\widetilde{H}_0(X)=0.
}
\]

For spheres,

\[
\boxed{
\widetilde{H}_k(S^n)
\cong
\begin{cases}
\Z,
&
k=n,
\\
0,
&
k\neq n.
\end{cases}
}
\]

Reduced homology often makes formulas cleaner.

%%%%%%%%%%%%%%%%%%%%%%%%%%%%%%%%%%%%%%%%%%%%%%%%%%%%%%%%%%%%
\subsection{Homology with Coefficients}
\label{subsec:homology-coefficients}
%%%%%%%%%%%%%%%%%%%%%%%%%%%%%%%%%%%%%%%%%%%%%%%%%%%%%%%%%%%%

Homology need not use only integer coefficients.

One may construct

\[
\boxed{
H_n(X;G),
}
\]

where \(G\) is an abelian coefficient group.

Common choices include

\[
\boxed{
\Z,
\qquad
\Q,
\qquad
\Z/2\Z.
}
\]

The choice of coefficients can change what topological information is
visible.

%%%%%%%%%%%%%%%%%%%%%%%%%%%%%%%%%%%%%%%%%%%%%%%%%%%%%%%%%%%%
\subsection{Betti Numbers}
\label{subsec:betti-numbers}
%%%%%%%%%%%%%%%%%%%%%%%%%%%%%%%%%%%%%%%%%%%%%%%%%%%%%%%%%%%%

When homology is considered over a field such as \(\Q\), its dimensions are
called Betti numbers.

\begin{definitionbox}{Betti Number}
The \(n\)-th Betti number is

\[
\boxed{
b_n
=
\dim_{\Q}H_n(X;\Q).
}
\]
\end{definitionbox}

The number \(b_n\) roughly counts independent \(n\)-dimensional holes.

%%%%%%%%%%%%%%%%%%%%%%%%%%%%%%%%%%%%%%%%%%%%%%%%%%%%%%%%%%%%
\subsection{Examples of Betti Numbers}
\label{subsec:betti-number-examples}
%%%%%%%%%%%%%%%%%%%%%%%%%%%%%%%%%%%%%%%%%%%%%%%%%%%%%%%%%%%%

For the circle,

\[
\boxed{
b_0=1,
\qquad
b_1=1.
}
\]

For the torus,

\[
\boxed{
b_0=1,
\qquad
b_1=2,
\qquad
b_2=1.
}
\]

For the sphere \(S^2\),

\[
\boxed{
b_0=1,
\qquad
b_1=0,
\qquad
b_2=1.
}
\]

%%%%%%%%%%%%%%%%%%%%%%%%%%%%%%%%%%%%%%%%%%%%%%%%%%%%%%%%%%%%
\subsection{Torsion in Homology}
\label{subsec:torsion-homology}
%%%%%%%%%%%%%%%%%%%%%%%%%%%%%%%%%%%%%%%%%%%%%%%%%%%%%%%%%%%%

Homology groups may contain finite-order elements.

For example, a group may have the form

\[
\boxed{
H_n(X;\Z)
\cong
\Z^r
\oplus
\Z/m_1\Z
\oplus\cdots\oplus
\Z/m_k\Z.
}
\]

The free rank \(r\) contributes to the Betti number.

The finite cyclic factors represent \emph{torsion}.

Torsion contains topological information that Betti numbers alone do not
detect.

%%%%%%%%%%%%%%%%%%%%%%%%%%%%%%%%%%%%%%%%%%%%%%%%%%%%%%%%%%%%
\subsection{The Real Projective Plane}
\label{subsec:homology-real-projective-plane}
%%%%%%%%%%%%%%%%%%%%%%%%%%%%%%%%%%%%%%%%%%%%%%%%%%%%%%%%%%%%

The real projective plane

\[
\mathbb{RP}^2
\]

provides a standard example of torsion.

Its integral homology is

\[
\boxed{
H_0(\mathbb{RP}^2;\Z)\cong\Z,
}
\]

\[
\boxed{
H_1(\mathbb{RP}^2;\Z)
\cong
\Z/2\Z,
}
\]

and

\[
\boxed{
H_2(\mathbb{RP}^2;\Z)=0.
}
\]

The factor

\[
\Z/2\Z
\]

cannot be detected merely by the first Betti number over \(\Q\).

%%%%%%%%%%%%%%%%%%%%%%%%%%%%%%%%%%%%%%%%%%%%%%%%%%%%%%%%%%%%
\subsection{Simplicial Homology}
\label{subsec:simplicial-homology}
%%%%%%%%%%%%%%%%%%%%%%%%%%%%%%%%%%%%%%%%%%%%%%%%%%%%%%%%%%%%

When a space is represented by a simplicial complex, its homology can be
computed using the simplicial chain complex.

This procedure is finite and algebraic when the complex has finitely many
simplices.

The main steps are

\[
\boxed{
\text{simplices}
\longrightarrow
\text{chain groups}
\longrightarrow
\text{boundary maps}
\longrightarrow
\text{kernels and images}
\longrightarrow
\text{homology}.
}
\]

%%%%%%%%%%%%%%%%%%%%%%%%%%%%%%%%%%%%%%%%%%%%%%%%%%%%%%%%%%%%
\subsection{Boundary Matrices}
\label{subsec:boundary-matrices}
%%%%%%%%%%%%%%%%%%%%%%%%%%%%%%%%%%%%%%%%%%%%%%%%%%%%%%%%%%%%

Once bases for the chain groups are chosen, each boundary homomorphism can
be represented by a matrix.

Thus

\[
\partial_n
\]

becomes a matrix

\[
D_n.
\]

Then

\[
\boxed{
H_n
=
\frac{
\ker D_n
}{
\operatorname{im}D_{n+1}
}.
}
\]

This creates a direct connection between topology and linear algebra.

%%%%%%%%%%%%%%%%%%%%%%%%%%%%%%%%%%%%%%%%%%%%%%%%%%%%%%%%%%%%
\subsection{Homology Over a Field}
\label{subsec:homology-over-field}
%%%%%%%%%%%%%%%%%%%%%%%%%%%%%%%%%%%%%%%%%%%%%%%%%%%%%%%%%%%%

If coefficients are taken in a field \(F\), the chain groups become vector
spaces.

Then

\[
\boxed{
\dim H_n
=
\dim\ker\partial_n
-
\dim\operatorname{im}\partial_{n+1}.
}
\]

By rank-nullity,

\[
\dim\ker\partial_n
=
\dim C_n
-
\operatorname{rank}(\partial_n).
\]

Therefore

\[
\boxed{
b_n
=
\dim C_n
-
\operatorname{rank}(\partial_n)
-
\operatorname{rank}(\partial_{n+1})
}
\]

when the coefficient field and finite-dimensional hypotheses are understood.

%%%%%%%%%%%%%%%%%%%%%%%%%%%%%%%%%%%%%%%%%%%%%%%%%%%%%%%%%%%%
\subsection{Worked Example: Homology of a Triangle Boundary}
\label{subsec:worked-homology-triangle-boundary}
%%%%%%%%%%%%%%%%%%%%%%%%%%%%%%%%%%%%%%%%%%%%%%%%%%%%%%%%%%%%

\begin{workedexamplebox}{A Triangular Loop}

Consider a simplicial complex consisting of three vertices

\[
v_0,v_1,v_2
\]

and the three edges

\[
e_{01},
\qquad
e_{12},
\qquad
e_{20},
\]

but with no filled \(2\)-simplex.

The chain

\[
c=e_{01}+e_{12}+e_{20}
\]

forms a cycle because

\[
\partial_1 c=0.
\]

Since there are no \(2\)-simplices,

\[
C_2=0,
\]

so

\[
B_1=\operatorname{im}\partial_2=0.
\]

Therefore the loop represents a nonzero homology class.

\begin{answerbox}
\[
\boxed{
H_1\cong\Z.
}
\]

The unfilled triangular loop has one independent one-dimensional hole.
\end{answerbox}
\end{workedexamplebox}

%%%%%%%%%%%%%%%%%%%%%%%%%%%%%%%%%%%%%%%%%%%%%%%%%%%%%%%%%%%%
\subsection{Filled Versus Unfilled Triangle}
\label{subsec:filled-unfilled-triangle}
%%%%%%%%%%%%%%%%%%%%%%%%%%%%%%%%%%%%%%%%%%%%%%%%%%%%%%%%%%%%

Now add the filled triangle

\[
[v_0,v_1,v_2].
\]

Its boundary is exactly the surrounding edge cycle.

Therefore the cycle becomes a boundary:

\[
\boxed{
c=\partial_2[v_0,v_1,v_2].
}
\]

Hence its homology class becomes zero.

This demonstrates the fundamental principle:

\[
\boxed{
\text{a cycle represents a hole only when it does not bound}.
}
\]

%%%%%%%%%%%%%%%%%%%%%%%%%%%%%%%%%%%%%%%%%%%%%%%%%%%%%%%%%%%%
\subsection{Singular Simplices}
\label{subsec:singular-simplices}
%%%%%%%%%%%%%%%%%%%%%%%%%%%%%%%%%%%%%%%%%%%%%%%%%%%%%%%%%%%%

Simplicial homology requires a simplicial representation of the space.

Singular homology avoids this restriction.

A \emph{singular \(n\)-simplex} in \(X\) is a continuous map

\[
\boxed{
\sigma:\Delta^n\to X.
}
\]

The simplex may bend, fold, or map noninjectively into the space.

%%%%%%%%%%%%%%%%%%%%%%%%%%%%%%%%%%%%%%%%%%%%%%%%%%%%%%%%%%%%
\subsection{Singular Chain Groups}
\label{subsec:singular-chain-groups}
%%%%%%%%%%%%%%%%%%%%%%%%%%%%%%%%%%%%%%%%%%%%%%%%%%%%%%%%%%%%

The singular chain group

\[
\boxed{
C_n(X)
}
\]

is the free abelian group generated by all singular \(n\)-simplices

\[
\sigma:\Delta^n\to X.
\]

Boundary maps are defined by restricting each singular simplex to the faces
of the standard simplex with alternating signs.

The resulting homology is called \emph{singular homology}.

%%%%%%%%%%%%%%%%%%%%%%%%%%%%%%%%%%%%%%%%%%%%%%%%%%%%%%%%%%%%
\subsection{Why Singular Homology Is Important}
\label{subsec:why-singular-homology}
%%%%%%%%%%%%%%%%%%%%%%%%%%%%%%%%%%%%%%%%%%%%%%%%%%%%%%%%%%%%

Singular homology can be defined for essentially every topological space
without first choosing a triangulation.

It is therefore one of the standard general homology theories.

For spaces admitting suitable triangulations, simplicial and singular
homology agree.

%%%%%%%%%%%%%%%%%%%%%%%%%%%%%%%%%%%%%%%%%%%%%%%%%%%%%%%%%%%%
\subsection{Homology Is Functorial}
\label{subsec:homology-functorial}
%%%%%%%%%%%%%%%%%%%%%%%%%%%%%%%%%%%%%%%%%%%%%%%%%%%%%%%%%%%%

A continuous map

\[
f:X\to Y
\]

induces homomorphisms

\[
\boxed{
f_*:
H_n(X)
\to
H_n(Y).
}
\]

These satisfy

\[
\boxed{
(\operatorname{id}_X)_*
=
\operatorname{id}_{H_n(X)}
}
\]

and

\[
\boxed{
(g\circ f)_*
=
g_*\circ f_*.
}
\]

Thus homology, like the fundamental group, behaves functorially.

%%%%%%%%%%%%%%%%%%%%%%%%%%%%%%%%%%%%%%%%%%%%%%%%%%%%%%%%%%%%
\subsection{Homotopy Invariance of Homology}
\label{subsec:homology-homotopy-invariance}
%%%%%%%%%%%%%%%%%%%%%%%%%%%%%%%%%%%%%%%%%%%%%%%%%%%%%%%%%%%%

If

\[
f,g:X\to Y
\]

are homotopic, then they induce the same map on homology:

\[
\boxed{
f_*=g_*.
}
\]

Consequently,

\[
\boxed{
X\simeq Y
\Longrightarrow
H_n(X)\cong H_n(Y)
}
\]

for every \(n\).

Thus homology is a homotopy invariant.

%%%%%%%%%%%%%%%%%%%%%%%%%%%%%%%%%%%%%%%%%%%%%%%%%%%%%%%%%%%%
\subsection{Euler Characteristic}
\label{subsec:euler-characteristic-topology}
%%%%%%%%%%%%%%%%%%%%%%%%%%%%%%%%%%%%%%%%%%%%%%%%%%%%%%%%%%%%

For a finite simplicial complex, let

\[
f_k
\]

denote the number of \(k\)-simplices.

The Euler characteristic is

\[
\boxed{
\chi(X)
=
\sum_{k\geq0}
(-1)^k f_k.
}
\]

The lowercase Greek letter

\[
\chi
\]

is read ``chi.''

For a two-dimensional complex,

\[
\boxed{
\chi
=
V-E+F.
}
\]

%%%%%%%%%%%%%%%%%%%%%%%%%%%%%%%%%%%%%%%%%%%%%%%%%%%%%%%%%%%%
\subsection{Euler Characteristic from Homology}
\label{subsec:euler-characteristic-homology}
%%%%%%%%%%%%%%%%%%%%%%%%%%%%%%%%%%%%%%%%%%%%%%%%%%%%%%%%%%%%

Under suitable finite-dimensional hypotheses,

\[
\boxed{
\chi(X)
=
\sum_{k\geq0}
(-1)^k b_k.
}
\]

Thus Euler characteristic can be computed either from cells or simplices,
or from homology.

This is a striking connection between combinatorics and topology.

%%%%%%%%%%%%%%%%%%%%%%%%%%%%%%%%%%%%%%%%%%%%%%%%%%%%%%%%%%%%
\subsection{Euler Characteristic Examples}
\label{subsec:euler-characteristic-examples}
%%%%%%%%%%%%%%%%%%%%%%%%%%%%%%%%%%%%%%%%%%%%%%%%%%%%%%%%%%%%

For a point,

\[
\boxed{
\chi=1.
}
\]

For a circle,

\[
\boxed{
\chi(S^1)=0.
}
\]

For a sphere,

\[
\boxed{
\chi(S^2)=2.
}
\]

For a torus,

\[
\boxed{
\chi(T^2)=0.
}
\]

For a closed orientable surface of genus \(g\),

\[
\boxed{
\chi(\Sigma_g)
=
2-2g.
}
\]

%%%%%%%%%%%%%%%%%%%%%%%%%%%%%%%%%%%%%%%%%%%%%%%%%%%%%%%%%%%%
\subsection{Exact Sequences}
\label{subsec:exact-sequences-topology}
%%%%%%%%%%%%%%%%%%%%%%%%%%%%%%%%%%%%%%%%%%%%%%%%%%%%%%%%%%%%

A sequence of homomorphisms

\[
\cdots
\longrightarrow
A_{n+1}
\xrightarrow{f_{n+1}}
A_n
\xrightarrow{f_n}
A_{n-1}
\longrightarrow
\cdots
\]

is \emph{exact} at \(A_n\) when

\[
\boxed{
\operatorname{im}(f_{n+1})
=
\ker(f_n).
}
\]

An exact sequence encodes the idea that everything mapped to zero at one
stage comes from the previous stage.

%%%%%%%%%%%%%%%%%%%%%%%%%%%%%%%%%%%%%%%%%%%%%%%%%%%%%%%%%%%%
\subsection{Short Exact Sequences}
\label{subsec:short-exact-sequences-topology}
%%%%%%%%%%%%%%%%%%%%%%%%%%%%%%%%%%%%%%%%%%%%%%%%%%%%%%%%%%%%

A short exact sequence has the form

\[
\boxed{
0
\longrightarrow
A
\xrightarrow{f}
B
\xrightarrow{g}
C
\longrightarrow
0.
}
\]

Exactness implies:

\[
f
\]

is injective,

\[
g
\]

is surjective, and

\[
\boxed{
\operatorname{im}(f)
=
\ker(g).
}
\]

Exact sequences are among the central organizational tools of algebraic
topology.

%%%%%%%%%%%%%%%%%%%%%%%%%%%%%%%%%%%%%%%%%%%%%%%%%%%%%%%%%%%%
\subsection{Long Exact Sequences}
\label{subsec:long-exact-sequences-topology}
%%%%%%%%%%%%%%%%%%%%%%%%%%%%%%%%%%%%%%%%%%%%%%%%%%%%%%%%%%%%

Topological constructions often produce long exact sequences such as

\[
\cdots
\to
H_n(A)
\to
H_n(X)
\to
H_n(X,A)
\to
H_{n-1}(A)
\to
\cdots.
\]

These sequences allow unknown homology groups to be determined from known
ones.

%%%%%%%%%%%%%%%%%%%%%%%%%%%%%%%%%%%%%%%%%%%%%%%%%%%%%%%%%%%%
\subsection{Relative Homology}
\label{subsec:relative-homology}
%%%%%%%%%%%%%%%%%%%%%%%%%%%%%%%%%%%%%%%%%%%%%%%%%%%%%%%%%%%%

Given

\[
A\subseteq X,
\]

relative homology studies \(X\) while treating \(A\) as collapsed or
algebraically negligible.

The relative chain groups are

\[
\boxed{
C_n(X,A)
=
C_n(X)/C_n(A).
}
\]

Their homology groups are denoted

\[
\boxed{
H_n(X,A).
}
\]

Relative homology is essential in excision, manifold theory, and many
computations.

%%%%%%%%%%%%%%%%%%%%%%%%%%%%%%%%%%%%%%%%%%%%%%%%%%%%%%%%%%%%
\subsection{Mayer--Vietoris Sequence}
\label{subsec:mayer-vietoris}
%%%%%%%%%%%%%%%%%%%%%%%%%%%%%%%%%%%%%%%%%%%%%%%%%%%%%%%%%%%%

Suppose a space is decomposed into overlapping pieces

\[
X=A\cup B.
\]

The Mayer--Vietoris sequence relates the homology of

\[
A,
\qquad
B,
\qquad
A\cap B,
\qquad
X.
\]

A portion has the form

\[
\boxed{
\cdots
\to
H_n(A\cap B)
\to
H_n(A)\oplus H_n(B)
\to
H_n(X)
\to
H_{n-1}(A\cap B)
\to
\cdots.
}
\]

This is the homological analogue of decomposing a space into simpler pieces.

%%%%%%%%%%%%%%%%%%%%%%%%%%%%%%%%%%%%%%%%%%%%%%%%%%%%%%%%%%%%
\subsection{Excision}
\label{subsec:excision-topology}
%%%%%%%%%%%%%%%%%%%%%%%%%%%%%%%%%%%%%%%%%%%%%%%%%%%%%%%%%%%%

The Excision Theorem states, roughly, that under appropriate conditions one
may remove a suitably contained portion of a pair

\[
(X,A)
\]

without changing its relative homology.

This allows local pieces of spaces to be removed when they do not affect the
relevant relative topology.

Excision is one of the foundational computational principles of singular
homology.

%%%%%%%%%%%%%%%%%%%%%%%%%%%%%%%%%%%%%%%%%%%%%%%%%%%%%%%%%%%%
\subsection{Cell Complexes}
\label{subsec:cell-complexes}
%%%%%%%%%%%%%%%%%%%%%%%%%%%%%%%%%%%%%%%%%%%%%%%%%%%%%%%%%%%%

Not every useful decomposition must use simplices.

A \emph{cell complex} builds a space from disks of various dimensions.

A \(0\)-cell is a point.

A \(1\)-cell behaves like an interval.

A \(2\)-cell behaves like a disk.

An \(n\)-cell behaves like an open \(n\)-dimensional ball.

%%%%%%%%%%%%%%%%%%%%%%%%%%%%%%%%%%%%%%%%%%%%%%%%%%%%%%%%%%%%
\subsection{CW Complexes}
\label{subsec:cw-complexes}
%%%%%%%%%%%%%%%%%%%%%%%%%%%%%%%%%%%%%%%%%%%%%%%%%%%%%%%%%%%%

A CW complex is built inductively by attaching cells.

The process begins with a collection of \(0\)-cells.

Then \(1\)-cells are attached.

Next \(2\)-cells are attached.

The process continues dimension by dimension.

Many important spaces in algebraic topology admit useful CW structures.

%%%%%%%%%%%%%%%%%%%%%%%%%%%%%%%%%%%%%%%%%%%%%%%%%%%%%%%%%%%%
\subsection{Example: CW Structure of the Circle}
\label{subsec:cw-circle}
%%%%%%%%%%%%%%%%%%%%%%%%%%%%%%%%%%%%%%%%%%%%%%%%%%%%%%%%%%%%

The circle can be built from

\[
\boxed{
\text{one }0\text{-cell}
}
\]

and

\[
\boxed{
\text{one }1\text{-cell}.
}
\]

The endpoints of the \(1\)-cell are attached to the same \(0\)-cell.

Therefore

\[
\chi(S^1)
=
1-1
=
0.
\]

%%%%%%%%%%%%%%%%%%%%%%%%%%%%%%%%%%%%%%%%%%%%%%%%%%%%%%%%%%%%
\subsection{Example: CW Structure of the Torus}
\label{subsec:cw-torus}
%%%%%%%%%%%%%%%%%%%%%%%%%%%%%%%%%%%%%%%%%%%%%%%%%%%%%%%%%%%%

The torus admits a CW decomposition with

\[
\boxed{
1\text{ zero-cell},
}
\]

\[
\boxed{
2\text{ one-cells},
}
\]

and

\[
\boxed{
1\text{ two-cell}.
}
\]

Therefore

\[
\boxed{
\chi(T^2)
=
1-2+1
=
0.
}
\]

This decomposition also reflects the two fundamental loop directions of the
torus.

%%%%%%%%%%%%%%%%%%%%%%%%%%%%%%%%%%%%%%%%%%%%%%%%%%%%%%%%%%%%
\subsection{Cohomology}
\label{subsec:cohomology-introduction}
%%%%%%%%%%%%%%%%%%%%%%%%%%%%%%%%%%%%%%%%%%%%%%%%%%%%%%%%%%%%

Homology is constructed from chain groups and boundary maps.

Cohomology reverses the algebraic direction.

Given a chain group \(C_n\) and coefficient group \(G\), define the
\(n\)-cochain group by

\[
\boxed{
C^n(X;G)
=
\operatorname{Hom}(C_n(X),G).
}
\]

The notation

\[
\operatorname{Hom}
\]

denotes the set or group of homomorphisms.

%%%%%%%%%%%%%%%%%%%%%%%%%%%%%%%%%%%%%%%%%%%%%%%%%%%%%%%%%%%%
\subsection{Coboundary Operator}
\label{subsec:coboundary-operator}
%%%%%%%%%%%%%%%%%%%%%%%%%%%%%%%%%%%%%%%%%%%%%%%%%%%%%%%%%%%%

The boundary maps induce coboundary maps

\[
\boxed{
\delta^n:
C^n
\to
C^{n+1}.
}
\]

The lowercase Greek letter

\[
\delta
\]

is read ``delta.''

These satisfy

\[
\boxed{
\delta^{n+1}\circ\delta^n=0.
}
\]

Thus cochains form a cochain complex.

%%%%%%%%%%%%%%%%%%%%%%%%%%%%%%%%%%%%%%%%%%%%%%%%%%%%%%%%%%%%
\subsection{Cocycles and Coboundaries}
\label{subsec:cocycles-coboundaries}
%%%%%%%%%%%%%%%%%%%%%%%%%%%%%%%%%%%%%%%%%%%%%%%%%%%%%%%%%%%%

The group of \(n\)-cocycles is

\[
\boxed{
Z^n
=
\ker(\delta^n),
}
\]

while the group of \(n\)-coboundaries is

\[
\boxed{
B^n
=
\operatorname{im}(\delta^{n-1}).
}
\]

Since

\[
\delta^2=0,
\]

we have

\[
\boxed{
B^n\subseteq Z^n.
}
\]

%%%%%%%%%%%%%%%%%%%%%%%%%%%%%%%%%%%%%%%%%%%%%%%%%%%%%%%%%%%%
\subsection{Cohomology Groups}
\label{subsec:cohomology-groups}
%%%%%%%%%%%%%%%%%%%%%%%%%%%%%%%%%%%%%%%%%%%%%%%%%%%%%%%%%%%%

\begin{definitionbox}{Cohomology Group}
The \(n\)-th cohomology group is

\[
\boxed{
H^n(X;G)
=
\frac{
Z^n(X;G)
}{
B^n(X;G)
}.
}
\]
\end{definitionbox}

Homology and cohomology contain closely related information, but cohomology
has additional algebraic structure.

%%%%%%%%%%%%%%%%%%%%%%%%%%%%%%%%%%%%%%%%%%%%%%%%%%%%%%%%%%%%
\subsection{Cup Product}
\label{subsec:cup-product}
%%%%%%%%%%%%%%%%%%%%%%%%%%%%%%%%%%%%%%%%%%%%%%%%%%%%%%%%%%%%

Cohomology groups admit a product called the \emph{cup product}:

\[
\boxed{
\smile:
H^p(X;R)\times H^q(X;R)
\to
H^{p+q}(X;R).
}
\]

The symbol

\[
\smile
\]

is read ``cup.''

Thus the direct sum

\[
\boxed{
H^*(X;R)
=
\bigoplus_{n\geq0}H^n(X;R)
}
\]

can be given the structure of a graded ring.

%%%%%%%%%%%%%%%%%%%%%%%%%%%%%%%%%%%%%%%%%%%%%%%%%%%%%%%%%%%%
\subsection{Cohomology Ring}
\label{subsec:cohomology-ring}
%%%%%%%%%%%%%%%%%%%%%%%%%%%%%%%%%%%%%%%%%%%%%%%%%%%%%%%%%%%%

The graded ring

\[
\boxed{
H^*(X;R)
}
\]

can distinguish spaces whose homology groups alone appear similar.

Thus cohomology does more than count holes.

It also records how cohomology classes interact multiplicatively.

%%%%%%%%%%%%%%%%%%%%%%%%%%%%%%%%%%%%%%%%%%%%%%%%%%%%%%%%%%%%
\subsection{Universal Coefficient Philosophy}
\label{subsec:universal-coefficient-philosophy}
%%%%%%%%%%%%%%%%%%%%%%%%%%%%%%%%%%%%%%%%%%%%%%%%%%%%%%%%%%%%

Homology and cohomology with different coefficient groups are related by
universal coefficient theorems.

These theorems explain how information such as

\[
\operatorname{Hom}
\]

and

\[
\operatorname{Ext}
\]

enters when changing between homology and cohomology.

The full algebraic statement is deferred until the necessary homological
algebra has been developed.

%%%%%%%%%%%%%%%%%%%%%%%%%%%%%%%%%%%%%%%%%%%%%%%%%%%%%%%%%%%%
\subsection{Degree of a Map}
\label{subsec:degree-of-map}
%%%%%%%%%%%%%%%%%%%%%%%%%%%%%%%%%%%%%%%%%%%%%%%%%%%%%%%%%%%%

A continuous map

\[
f:S^n\to S^n
\]

induces

\[
f_*:
H_n(S^n)
\to
H_n(S^n).
\]

Since

\[
H_n(S^n)\cong\Z,
\]

the homomorphism must have the form

\[
\boxed{
f_*(m)=d\,m
}
\]

for some integer \(d\).

This integer is called the \emph{degree} of \(f\):

\[
\boxed{
\deg(f)=d.
}
\]

%%%%%%%%%%%%%%%%%%%%%%%%%%%%%%%%%%%%%%%%%%%%%%%%%%%%%%%%%%%%
\subsection{Meaning of Degree}
\label{subsec:meaning-degree}
%%%%%%%%%%%%%%%%%%%%%%%%%%%%%%%%%%%%%%%%%%%%%%%%%%%%%%%%%%%%

The degree measures, with orientation, how many times the domain sphere
wraps around the target sphere.

For maps

\[
S^1\to S^1,
\]

degree agrees with the winding-number intuition.

Degree is an important invariant in topology, geometry, and differential
equations.

%%%%%%%%%%%%%%%%%%%%%%%%%%%%%%%%%%%%%%%%%%%%%%%%%%%%%%%%%%%%
\subsection{Brouwer Fixed-Point Theorem}
\label{subsec:brouwer-fixed-point}
%%%%%%%%%%%%%%%%%%%%%%%%%%%%%%%%%%%%%%%%%%%%%%%%%%%%%%%%%%%%

Algebraic topology can prove existence theorems that initially appear
analytic or geometric.

\begin{theorem}[Brouwer Fixed-Point Theorem]
Every continuous map

\[
f:D^n\to D^n
\]

has a fixed point.

That is, there exists

\[
x\in D^n
\]

such that

\[
\boxed{
f(x)=x.
}
\]
\end{theorem}

The theorem has many proofs using algebraic-topological ideas.

%%%%%%%%%%%%%%%%%%%%%%%%%%%%%%%%%%%%%%%%%%%%%%%%%%%%%%%%%%%%
\subsection{No-Retraction Theorem}
\label{subsec:no-retraction-theorem}
%%%%%%%%%%%%%%%%%%%%%%%%%%%%%%%%%%%%%%%%%%%%%%%%%%%%%%%%%%%%

A closely related result states:

\begin{theorem}[No-Retraction Theorem]
There is no continuous retraction

\[
\boxed{
r:D^n\to S^{n-1}
}
\]

from the disk onto its boundary sphere.
\end{theorem}

Homology provides a natural proof because the disk and sphere have different
homological behavior.

%%%%%%%%%%%%%%%%%%%%%%%%%%%%%%%%%%%%%%%%%%%%%%%%%%%%%%%%%%%%
\subsection{Borsuk--Ulam Theorem}
\label{subsec:borsuk-ulam}
%%%%%%%%%%%%%%%%%%%%%%%%%%%%%%%%%%%%%%%%%%%%%%%%%%%%%%%%%%%%

Another major theorem is:

\begin{theorem}[Borsuk--Ulam Theorem]
For every continuous map

\[
f:S^n\to\R^n,
\]

there exists

\[
x\in S^n
\]

such that

\[
\boxed{
f(x)=f(-x).
}
\]
\end{theorem}

Thus some pair of antipodal points must have the same image.

This theorem has applications in combinatorics, geometry, and economics.

%%%%%%%%%%%%%%%%%%%%%%%%%%%%%%%%%%%%%%%%%%%%%%%%%%%%%%%%%%%%
\subsection{Fundamental Group Versus First Homology}
\label{subsec:pi1-vs-h1}
%%%%%%%%%%%%%%%%%%%%%%%%%%%%%%%%%%%%%%%%%%%%%%%%%%%%%%%%%%%%

The fundamental group may be nonabelian.

Homology groups are always abelian.

There is a fundamental relationship:

\[
\boxed{
H_1(X;\Z)
\cong
\pi_1(X)_{\mathrm{ab}}
}
\]

for path-connected spaces.

Here

\[
\pi_1(X)_{\mathrm{ab}}
\]

denotes the abelianization of the fundamental group.

Thus first homology retains the abelian part of the information contained in
the fundamental group.

%%%%%%%%%%%%%%%%%%%%%%%%%%%%%%%%%%%%%%%%%%%%%%%%%%%%%%%%%%%%
\subsection{Example: Torus and Wedge of Two Circles}
\label{subsec:torus-wedge-comparison}
%%%%%%%%%%%%%%%%%%%%%%%%%%%%%%%%%%%%%%%%%%%%%%%%%%%%%%%%%%%%

Consider

\[
T^2
\]

and

\[
S^1\vee S^1.
\]

Their first homology groups are both

\[
\boxed{
H_1\cong\Z^2.
}
\]

However,

\[
\pi_1(T^2)
\cong
\Z^2
\]

is abelian, while

\[
\pi_1(S^1\vee S^1)
\cong
F_2
\]

is nonabelian.

Therefore the fundamental group distinguishes information that first
homology alone does not.

%%%%%%%%%%%%%%%%%%%%%%%%%%%%%%%%%%%%%%%%%%%%%%%%%%%%%%%%%%%%
\subsection{Homology Does Not Classify All Spaces}
\label{subsec:homology-not-complete}
%%%%%%%%%%%%%%%%%%%%%%%%%%%%%%%%%%%%%%%%%%%%%%%%%%%%%%%%%%%%

Having isomorphic homology groups does not imply that two spaces are
homeomorphic or homotopy equivalent.

Likewise, having isomorphic fundamental groups does not classify all spaces.

Algebraic topology therefore develops many different invariants.

Each invariant captures certain features while forgetting others.

\begin{conceptbox}
Topological invariants are tools for distinguishing spaces, not generally
complete descriptions of spaces.
\end{conceptbox}

%%%%%%%%%%%%%%%%%%%%%%%%%%%%%%%%%%%%%%%%%%%%%%%%%%%%%%%%%%%%
\subsection{Persistent Homology}
\label{subsec:persistent-homology-intro}
%%%%%%%%%%%%%%%%%%%%%%%%%%%%%%%%%%%%%%%%%%%%%%%%%%%%%%%%%%%%

A modern application of algebraic topology is \emph{persistent homology}.

Instead of studying one space, persistent homology studies a nested family

\[
\boxed{
X_0
\subseteq
X_1
\subseteq
X_2
\subseteq
\cdots.
}
\]

Homology classes may

\[
\boxed{
\text{appear},
\quad
\text{persist},
\quad
\text{disappear}
}
\]

as the parameter changes.

This is a central idea in topological data analysis.

%%%%%%%%%%%%%%%%%%%%%%%%%%%%%%%%%%%%%%%%%%%%%%%%%%%%%%%%%%%%
\subsection{Filtrations}
\label{subsec:topological-filtrations}
%%%%%%%%%%%%%%%%%%%%%%%%%%%%%%%%%%%%%%%%%%%%%%%%%%%%%%%%%%%%

\begin{definitionbox}{Filtration}
A \emph{filtration} is a nested family of spaces or complexes

\[
\boxed{
X_0
\subseteq
X_1
\subseteq
\cdots
\subseteq
X_m.
}
\]
\end{definitionbox}

Each inclusion induces homology maps

\[
H_n(X_i)
\to
H_n(X_j)
\]

for

\[
i\leq j.
\]

These maps track the lifetime of topological features.

%%%%%%%%%%%%%%%%%%%%%%%%%%%%%%%%%%%%%%%%%%%%%%%%%%%%%%%%%%%%
\subsection{Barcodes}
\label{subsec:persistence-barcodes}
%%%%%%%%%%%%%%%%%%%%%%%%%%%%%%%%%%%%%%%%%%%%%%%%%%%%%%%%%%%%

Persistent homology is often visualized using a \emph{barcode}.

Each interval represents the lifespan of a homological feature.

Long intervals are often interpreted as persistent structure.

Short intervals may correspond to small-scale structure or noise, depending
on the application.

The interpretation must be made carefully and depends on how the filtration
was constructed.

%%%%%%%%%%%%%%%%%%%%%%%%%%%%%%%%%%%%%%%%%%%%%%%%%%%%%%%%%%%%
\subsection{Algebraic Topology and Manifolds}
\label{subsec:algebraic-topology-manifolds}
%%%%%%%%%%%%%%%%%%%%%%%%%%%%%%%%%%%%%%%%%%%%%%%%%%%%%%%%%%%%

Manifolds have local Euclidean structure but may possess complicated global
topology.

Algebraic invariants help distinguish them.

For example, surfaces may be studied using

\[
\boxed{
\pi_1,
\quad
H_1,
\quad
H_2,
\quad
\chi.
}
\]

Later sections will combine these invariants with geometric and
differentiable structure.

%%%%%%%%%%%%%%%%%%%%%%%%%%%%%%%%%%%%%%%%%%%%%%%%%%%%%%%%%%%%
\subsection{Algebraic Topology and Knot Theory}
\label{subsec:algebraic-topology-knots}
%%%%%%%%%%%%%%%%%%%%%%%%%%%%%%%%%%%%%%%%%%%%%%%%%%%%%%%%%%%%

Given a knot

\[
K\subseteq S^3,
\]

one can study its complement

\[
\boxed{
S^3\setminus K.
}
\]

The fundamental group

\[
\boxed{
\pi_1(S^3\setminus K)
}
\]

is called the \emph{knot group}.

It provides an important invariant of knots.

This connection will be developed further in the section on knot theory.

%%%%%%%%%%%%%%%%%%%%%%%%%%%%%%%%%%%%%%%%%%%%%%%%%%%%%%%%%%%%
\subsection{Algebraic Topology and Differential Forms}
\label{subsec:algebraic-topology-differential-forms}
%%%%%%%%%%%%%%%%%%%%%%%%%%%%%%%%%%%%%%%%%%%%%%%%%%%%%%%%%%%%

On smooth manifolds, differential forms produce another cohomology theory:
\emph{de Rham cohomology}.

It is defined using the exterior derivative

\[
d.
\]

Closed forms satisfy

\[
d\omega=0,
\]

while exact forms have the form

\[
\omega=d\eta.
\]

Thus de Rham cohomology has the same algebraic pattern:

\[
\boxed{
\frac{\text{closed objects}}{\text{exact objects}}.
}
\]

%%%%%%%%%%%%%%%%%%%%%%%%%%%%%%%%%%%%%%%%%%%%%%%%%%%%%%%%%%%%
\subsection{The de Rham Theorem}
\label{subsec:de-rham-theorem-preview}
%%%%%%%%%%%%%%%%%%%%%%%%%%%%%%%%%%%%%%%%%%%%%%%%%%%%%%%%%%%%

A fundamental theorem connects differential geometry with algebraic
topology.

\begin{theorem}[de Rham Theorem]
For a smooth manifold \(M\), de Rham cohomology is naturally isomorphic to
singular cohomology with real coefficients:

\[
\boxed{
H_{\mathrm{dR}}^k(M)
\cong
H^k(M;\R).
}
\]
\end{theorem}

Thus differential forms can recover topological information.

%%%%%%%%%%%%%%%%%%%%%%%%%%%%%%%%%%%%%%%%%%%%%%%%%%%%%%%%%%%%
\subsection{Common Errors}
\label{subsec:algebraic-topology-common-errors}
%%%%%%%%%%%%%%%%%%%%%%%%%%%%%%%%%%%%%%%%%%%%%%%%%%%%%%%%%%%%

\subsubsection{Confusing Homotopy with Homeomorphism}

Homeomorphic spaces are homotopy equivalent.

The converse is false.

Homotopy equivalence is a weaker relation.

%%%%%%%%%%%%%%%%%%%%%%%%%%%%%%%%%%%%%%%%%%%%%%%%%%%%%%%%%%%%
\subsubsection{Assuming Homotopy Preserves Geometry}

A homotopy may dramatically change distances, angles, lengths, and
curvature.

It preserves homotopy-theoretic information, not metric geometry.

%%%%%%%%%%%%%%%%%%%%%%%%%%%%%%%%%%%%%%%%%%%%%%%%%%%%%%%%%%%%
\subsubsection{Ignoring Basepoints}

The fundamental group is formally

\[
\pi_1(X,x_0).
\]

In path-connected spaces different basepoints produce isomorphic groups, but
the basepoint should not simply be ignored when defining induced maps.

%%%%%%%%%%%%%%%%%%%%%%%%%%%%%%%%%%%%%%%%%%%%%%%%%%%%%%%%%%%%
\subsubsection{Assuming Every Loop Is Nontrivial}

A loop may look complicated while still being contractible.

The relevant question is its homotopy class.

%%%%%%%%%%%%%%%%%%%%%%%%%%%%%%%%%%%%%%%%%%%%%%%%%%%%%%%%%%%%
\subsubsection{Counting Holes Too Literally}

The statement that

\[
H_n
\]

``counts \(n\)-dimensional holes'' is useful intuition, but homology groups
contain algebraic structure that cannot always be reduced to a simple visual
count.

%%%%%%%%%%%%%%%%%%%%%%%%%%%%%%%%%%%%%%%%%%%%%%%%%%%%%%%%%%%%
\subsubsection{Confusing Cycles and Boundaries}

Every boundary is a cycle because

\[
\partial^2=0.
\]

But not every cycle is a boundary.

The difference between them is precisely what homology measures.

%%%%%%%%%%%%%%%%%%%%%%%%%%%%%%%%%%%%%%%%%%%%%%%%%%%%%%%%%%%%
\subsubsection{Forgetting Orientation Signs}

The boundary formula

\[
\partial[v_0,\ldots,v_n]
=
\sum_i(-1)^i
[v_0,\ldots,\widehat{v_i},\ldots,v_n]
\]

requires alternating signs.

Without them,

\[
\partial^2=0
\]

would fail.

%%%%%%%%%%%%%%%%%%%%%%%%%%%%%%%%%%%%%%%%%%%%%%%%%%%%%%%%%%%%
\subsubsection{Confusing Chain Groups with Homology Groups}

The chain group

\[
C_n
\]

contains all formal \(n\)-chains.

The homology group is

\[
\boxed{
H_n
=
\ker\partial_n/
\operatorname{im}\partial_{n+1}.
}
\]

They are different algebraic objects.

%%%%%%%%%%%%%%%%%%%%%%%%%%%%%%%%%%%%%%%%%%%%%%%%%%%%%%%%%%%%
\subsubsection{Assuming Homology Groups Are Always Free}

Integral homology can contain torsion.

For example,

\[
H_1(\mathbb{RP}^2;\Z)
\cong
\Z/2\Z.
\]

%%%%%%%%%%%%%%%%%%%%%%%%%%%%%%%%%%%%%%%%%%%%%%%%%%%%%%%%%%%%
\subsubsection{Confusing Betti Numbers with Full Homology}

Betti numbers record ranks over a field.

They do not record integral torsion.

%%%%%%%%%%%%%%%%%%%%%%%%%%%%%%%%%%%%%%%%%%%%%%%%%%%%%%%%%%%%
\subsubsection{Assuming Same Homology Means Same Space}

Different spaces can have identical homology groups.

Homology is an invariant, not a complete classification.

%%%%%%%%%%%%%%%%%%%%%%%%%%%%%%%%%%%%%%%%%%%%%%%%%%%%%%%%%%%%
\subsubsection{Assuming the Fundamental Group Is Abelian}

Fundamental groups can be nonabelian.

For example,

\[
\pi_1(S^1\vee S^1)
\cong
F_2.
\]

%%%%%%%%%%%%%%%%%%%%%%%%%%%%%%%%%%%%%%%%%%%%%%%%%%%%%%%%%%%%
\subsubsection{Confusing \texorpdfstring{\(\pi_1\)}{pi1} with
\texorpdfstring{\(H_1\)}{H1}}

The first homology group is the abelianization of the fundamental group under
standard path-connected hypotheses:

\[
H_1(X;\Z)
\cong
\pi_1(X)_{\mathrm{ab}}.
\]

Thus \(H_1\) generally contains less information.

%%%%%%%%%%%%%%%%%%%%%%%%%%%%%%%%%%%%%%%%%%%%%%%%%%%%%%%%%%%%
\subsubsection{Assuming Every Space Has a Universal Cover}

Existence of a universal cover requires hypotheses.

Path connectedness alone is not sufficient.

Standard sufficient assumptions include local path connectedness and
semilocal simple connectedness.

%%%%%%%%%%%%%%%%%%%%%%%%%%%%%%%%%%%%%%%%%%%%%%%%%%%%%%%%%%%%
\subsubsection{Confusing a Covering Space with a Quotient Space}

Covering spaces and quotient spaces are related in many examples but are
different constructions.

A covering map must satisfy a specific local homeomorphism condition.

%%%%%%%%%%%%%%%%%%%%%%%%%%%%%%%%%%%%%%%%%%%%%%%%%%%%%%%%%%%%
\subsubsection{Assuming Euler Characteristic Completely Determines Topology}

Many nonhomeomorphic spaces have the same Euler characteristic.

It is useful but relatively coarse.

%%%%%%%%%%%%%%%%%%%%%%%%%%%%%%%%%%%%%%%%%%%%%%%%%%%%%%%%%%%%
\subsubsection{Confusing Homology and Cohomology}

Homology is constructed from chains.

Cohomology is constructed from cochains.

They are related but not identical, and cohomology carries additional
multiplicative structure through the cup product.

%%%%%%%%%%%%%%%%%%%%%%%%%%%%%%%%%%%%%%%%%%%%%%%%%%%%%%%%%%%%
\subsection{Applications and Connections}
\label{subsec:algebraic-topology-applications}
%%%%%%%%%%%%%%%%%%%%%%%%%%%%%%%%%%%%%%%%%%%%%%%%%%%%%%%%%%%%

\begin{applicationbox}
\textbf{Group Theory.}
Fundamental groups translate loop structure into groups. Free groups,
presentations, quotient groups, and abelianization appear naturally.

\medskip

\textbf{Linear Algebra.}
Boundary operators can be represented by matrices, allowing homology over a
field to be computed using kernels, images, ranks, and nullities.

\medskip

\textbf{Homological Algebra.}
Chain complexes, exact sequences, homology, cohomology, \(\operatorname{Hom}\),
and later \(\operatorname{Ext}\) and \(\operatorname{Tor}\) provide the
algebraic machinery behind many topological constructions.

\medskip

\textbf{Geometry.}
Algebraic invariants distinguish geometric spaces even when visual intuition
is insufficient.

\medskip

\textbf{Differential Geometry.}
de Rham cohomology translates differential forms into topological
information.

\medskip

\textbf{Manifold Theory.}
Fundamental groups, homology, cohomology, degree, and orientation are central
tools for studying manifolds.

\medskip

\textbf{Knot Theory.}
The topology of knot complements produces invariants such as knot groups and
homology groups.

\medskip

\textbf{Fixed-Point Theory.}
Homology and degree methods lead to results such as the Brouwer Fixed-Point
Theorem.

\medskip

\textbf{Topological Data Analysis.}
Persistent homology tracks topological features across a filtration and
provides computable summaries of data shape.

\medskip

\textbf{Physics.}
Homotopy and cohomology appear in gauge theory, defects, field theory,
topological phases, and mathematical models of spacetime.
\end{applicationbox}

%%%%%%%%%%%%%%%%%%%%%%%%%%%%%%%%%%%%%%%%%%%%%%%%%%%%%%%%%%%%
\subsection{Exercises}
\label{subsec:algebraic-topology-exercises}
%%%%%%%%%%%%%%%%%%%%%%%%%%%%%%%%%%%%%%%%%%%%%%%%%%%%%%%%%%%%

\begin{exercise}[1]
What is the central philosophy of algebraic topology?
\end{exercise}

\begin{exercise}[1]
Define a loop based at \(x_0\).
\end{exercise}

\begin{exercise}[1]
Define the reverse of a path \(\gamma\).
\end{exercise}

\begin{exercise}[1]
Define path concatenation.
\end{exercise}

\begin{exercise}[1]
Define a homotopy between two maps.
\end{exercise}

\begin{exercise}[1]
What does

\[
f\simeq g
\]

mean?
\end{exercise}

\begin{exercise}[1]
Define homotopy equivalence.
\end{exercise}

\begin{exercise}[1]
Define a contractible space.
\end{exercise}

\begin{exercise}[1]
Give three examples of contractible spaces.
\end{exercise}

\begin{exercise}[1]
Define a retraction.
\end{exercise}

\begin{exercise}[1]
Define a deformation retract.
\end{exercise}

\begin{exercise}[1]
What does the fundamental group measure intuitively?
\end{exercise}

\begin{exercise}[1]
Compute

\[
\pi_1(\R^n).
\]
\end{exercise}

\begin{exercise}[1]
State

\[
\pi_1(S^1).
\]
\end{exercise}

\begin{exercise}[1]
State

\[
\pi_1(T^2).
\]
\end{exercise}

\begin{exercise}[1]
What does it mean for a space to be simply connected?
\end{exercise}

\begin{exercise}[1]
Give an example of a path-connected space that is not simply connected.
\end{exercise}

\begin{exercise}[1]
Define a covering map.
\end{exercise}

\begin{exercise}[1]
Give a standard universal covering map of \(S^1\).
\end{exercise}

\begin{exercise}[1]
What is a deck transformation?
\end{exercise}

\begin{exercise}[1]
What is an \(n\)-simplex?
\end{exercise}

\begin{exercise}[1]
Describe \(0\)-, \(1\)-, \(2\)-, and \(3\)-simplices.
\end{exercise}

\begin{exercise}[1]
Define a simplicial complex.
\end{exercise}

\begin{exercise}[1]
Define the chain group \(C_n\).
\end{exercise}

\begin{exercise}[1]
State the formula for the boundary of an oriented \(n\)-simplex.
\end{exercise}

\begin{exercise}[1]
What does

\[
\partial^2=0
\]

mean?
\end{exercise}

\begin{exercise}[1]
Define an \(n\)-cycle.
\end{exercise}

\begin{exercise}[1]
Define an \(n\)-boundary.
\end{exercise}

\begin{exercise}[1]
Define

\[
H_n.
\]
\end{exercise}

\begin{exercise}[1]
State the homology groups of a point.
\end{exercise}

\begin{exercise}[1]
State the homology groups of \(S^1\).
\end{exercise}

\begin{exercise}[1]
State the homology groups of \(S^2\).
\end{exercise}

\begin{exercise}[1]
What does \(H_0\) detect?
\end{exercise}

\begin{exercise}[1]
Define the \(n\)-th Betti number.
\end{exercise}

\begin{exercise}[1]
What is torsion in homology?
\end{exercise}

\begin{exercise}[1]
Define Euler characteristic.
\end{exercise}

\begin{exercise}[1]
Compute

\[
\chi(S^1),
\qquad
\chi(S^2),
\qquad
\chi(T^2).
\]
\end{exercise}

\begin{exercise}[2]
Show that path reversal satisfies

\[
\overline{\overline{\gamma}}=\gamma.
\]
\end{exercise}

\begin{exercise}[2]
Show that a convex subset of \(\R^n\) is contractible.
\end{exercise}

\begin{exercise}[2]
Construct an explicit deformation retraction

\[
\R^2\setminus\{0\}
\to
S^1.
\]
\end{exercise}

\begin{exercise}[2]
Use the previous exercise to determine

\[
\pi_1(\R^2\setminus\{0\}).
\]
\end{exercise}

\begin{exercise}[2]
Explain geometrically why

\[
\pi_1(S^1)\cong\Z.
\]
\end{exercise}

\begin{exercise}[2]
Why is

\[
S^1
\]

not simply connected?
\end{exercise}

\begin{exercise}[2]
Explain why

\[
S^2
\]

is simply connected.
\end{exercise}

\begin{exercise}[2]
Describe the two generators of

\[
\pi_1(T^2).
\]
\end{exercise}

\begin{exercise}[2]
Compare

\[
\pi_1(T^2)
\]

with

\[
\pi_1(S^1\vee S^1).
\]
\end{exercise}

\begin{exercise}[2]
For the covering

\[
p(t)=e^{2\pi i t},
\]

find the fiber over \(1\in S^1\).
\end{exercise}

\begin{exercise}[2]
Describe the deck transformations of

\[
\R\to S^1.
\]
\end{exercise}

\begin{exercise}[2]
Compute

\[
\partial[v_0,v_1].
\]
\end{exercise}

\begin{exercise}[2]
Compute

\[
\partial[v_0,v_1,v_2].
\]
\end{exercise}

\begin{exercise}[2]
Verify directly that

\[
\partial^2[v_0,v_1,v_2]=0.
\]
\end{exercise}

\begin{exercise}[2]
Explain why the boundary cycle of an unfilled triangle represents a
nontrivial first homology class.
\end{exercise}

\begin{exercise}[2]
Explain why filling the triangle makes the same cycle trivial in homology.
\end{exercise}

\begin{exercise}[2]
State the integral homology groups of \(T^2\).
\end{exercise}

\begin{exercise}[2]
Find the Betti numbers of \(T^2\).
\end{exercise}

\begin{exercise}[2]
Find the Betti numbers of \(S^2\).
\end{exercise}

\begin{exercise}[2]
Compute the Euler characteristic of a closed orientable surface of genus
\(4\).
\end{exercise}

\begin{exercise}[2]
Explain why

\[
H_1(\mathbb{RP}^2;\Q)=0
\]

does not mean that

\[
H_1(\mathbb{RP}^2;\Z)=0.
\]
\end{exercise}

\begin{exercise}[2]
What does it mean for a sequence of homomorphisms to be exact?
\end{exercise}

\begin{exercise}[2]
Explain the purpose of the Mayer--Vietoris sequence.
\end{exercise}

\begin{exercise}[2]
Describe a CW decomposition of \(S^1\).
\end{exercise}

\begin{exercise}[2]
Describe a CW decomposition of \(T^2\).
\end{exercise}

\begin{exercise}[2]
Define a cochain.
\end{exercise}

\begin{exercise}[2]
Define a cohomology group.
\end{exercise}

\begin{exercise}[2]
What additional structure does the cup product give to cohomology?
\end{exercise}

\begin{exercise}[3]
Prove that homotopy of maps is an equivalence relation.
\end{exercise}

\begin{exercise}[3]
Show that a deformation retract is a homotopy equivalence.
\end{exercise}

\begin{exercise}[3]
Prove that every contractible space has trivial fundamental group.
\end{exercise}

\begin{exercise}[3]
Show that a continuous based map

\[
f:(X,x_0)\to(Y,y_0)
\]

induces a homomorphism

\[
f_*:\pi_1(X,x_0)\to\pi_1(Y,y_0).
\]
\end{exercise}

\begin{exercise}[3]
Verify

\[
(g\circ f)_*=g_*\circ f_*
\]

for induced maps on fundamental groups.
\end{exercise}

\begin{exercise}[3]
Use fundamental groups to show that

\[
S^1
\]

is not homotopy equivalent to a point.
\end{exercise}

\begin{exercise}[3]
Use fundamental groups to show that

\[
T^2
\]

is not homotopy equivalent to \(S^2\).
\end{exercise}

\begin{exercise}[3]
Explain why the fundamental group of a wedge of two circles is nonabelian.
\end{exercise}

\begin{exercise}[3]
Use the Seifert--van Kampen theorem to compute the fundamental group of a
simple space assembled from overlapping pieces.
\end{exercise}

\begin{exercise}[3]
Prove that

\[
B_n\subseteq Z_n.
\]
\end{exercise}

\begin{exercise}[3]
Explain why the relation

\[
\partial^2=0
\]

is necessary for the quotient

\[
Z_n/B_n
\]

to be defined.
\end{exercise}

\begin{exercise}[3]
Construct boundary matrices for a simplicial circle consisting of three
vertices and three edges.
\end{exercise}

\begin{exercise}[3]
Use those matrices to compute

\[
H_0
\]

and

\[
H_1.
\]
\end{exercise}

\begin{exercise}[3]
Construct the boundary matrices for a filled triangle and compute its
homology.
\end{exercise}

\begin{exercise}[3]
Show that the Euler characteristic of \(S^2\) is \(2\) using the boundary of
a tetrahedron as a triangulation.
\end{exercise}

\begin{exercise}[3]
Show that

\[
\chi(\Sigma_g)=2-2g
\]

for a closed orientable surface using an appropriate cell decomposition.
\end{exercise}

\begin{exercise}[3]
Explain why homotopy-equivalent spaces have the same Euler characteristic
when the relevant finite-type hypotheses hold.
\end{exercise}

\begin{exercise}[3]
Explain why

\[
H_1(X;\Z)
\]

cannot detect the noncommutativity of

\[
\pi_1(X).
\]
\end{exercise}

\begin{exercise}[3]
Compare the first homology of \(T^2\) and \(S^1\vee S^1\).
\end{exercise}

\begin{exercise}[3]
Explain why the degree of a map

\[
S^n\to S^n
\]

must be an integer.
\end{exercise}

\begin{exercise}[4]
Prove

\[
\pi_1(S^1)\cong\Z
\]

using the covering map

\[
\R\to S^1.
\]
\end{exercise}

\begin{exercise}[4]
Study the Seifert--van Kampen theorem and derive a presentation of

\[
\pi_1(T^2).
\]
\end{exercise}

\begin{exercise}[4]
Compute the fundamental group of the real projective plane.
\end{exercise}

\begin{exercise}[4]
Investigate the fundamental group of the Klein bottle.
\end{exercise}

\begin{exercise}[4]
Research the classification of covering spaces by subgroups of the
fundamental group.
\end{exercise}

\begin{exercise}[4]
Compute the simplicial homology of the boundary of a tetrahedron.
\end{exercise}

\begin{exercise}[4]
Compute the cellular homology of the torus from its standard CW structure.
\end{exercise}

\begin{exercise}[4]
Compute the cellular homology of

\[
\mathbb{RP}^2.
\]
\end{exercise}

\begin{exercise}[4]
Use Mayer--Vietoris to compute the homology of \(S^1\).
\end{exercise}

\begin{exercise}[4]
Use Mayer--Vietoris to compute the homology of \(S^n\).
\end{exercise}

\begin{exercise}[4]
Research the Universal Coefficient Theorem for cohomology.
\end{exercise}

\begin{exercise}[4]
Research the K\"unneth theorem and explain how it relates the homology of a
product \(X\times Y\) to the homology of \(X\) and \(Y\).
\end{exercise}

\begin{exercise}[4]
Investigate the cohomology ring of the torus.
\end{exercise}

\begin{exercise}[4]
Research one algebraic-topological proof of the Brouwer Fixed-Point Theorem.
\end{exercise}

\begin{exercise}[4]
Research the Borsuk--Ulam theorem and give one application.
\end{exercise}

\begin{exercise}[4]
Construct a simple filtration of a planar point cloud and describe how
persistent \(H_0\) and \(H_1\) features might appear.
\end{exercise}

%%%%%%%%%%%%%%%%%%%%%%%%%%%%%%%%%%%%%%%%%%%%%%%%%%%%%%%%%%%%
\subsection{Conceptual Summary}
\label{subsec:algebraic-topology-summary}
%%%%%%%%%%%%%%%%%%%%%%%%%%%%%%%%%%%%%%%%%%%%%%%%%%%%%%%%%%%%

Algebraic topology assigns algebraic structures to topological spaces:

\[
\boxed{
\text{topology}
\longrightarrow
\text{algebra}.
}
\]

Homotopy formalizes continuous deformation:

\[
\boxed{
f\simeq g.
}
\]

Homotopy equivalence,

\[
\boxed{
X\simeq Y,
}
\]

is weaker than homeomorphism but preserves many important invariants.

The fundamental group

\[
\boxed{
\pi_1(X,x_0)
}
\]

studies homotopy classes of loops.

Important examples include

\[
\boxed{
\pi_1(S^1)\cong\Z,
}
\]

\[
\boxed{
\pi_1(T^2)\cong\Z^2,
}
\]

and

\[
\boxed{
\pi_1(S^1\vee S^1)\cong F_2.
}
\]

Covering spaces relate the geometry of a space to subgroups of its
fundamental group.

Homology begins with a chain complex

\[
\boxed{
\cdots
\to
C_{n+1}
\xrightarrow{\partial_{n+1}}
C_n
\xrightarrow{\partial_n}
C_{n-1}
\to
\cdots
}
\]

satisfying

\[
\boxed{
\partial^2=0.
}
\]

Cycles are

\[
\boxed{
Z_n=\ker\partial_n,
}
\]

boundaries are

\[
\boxed{
B_n=\operatorname{im}\partial_{n+1},
}
\]

and homology is

\[
\boxed{
H_n
=
Z_n/B_n.
}
\]

Betti numbers record the ranks of homology groups over a field.

Integral homology may additionally contain torsion.

Euler characteristic satisfies

\[
\boxed{
\chi(X)
=
\sum_n(-1)^n b_n
}
\]

under appropriate finite-type hypotheses.

Cohomology reverses the direction of the chain construction and produces

\[
\boxed{
H^n(X;G).
}
\]

The cup product gives

\[
\boxed{
H^*(X;R)
}
\]

the structure of a graded ring.

Algebraic topology therefore builds a bridge

\[
\boxed{
\text{geometry}
\longleftrightarrow
\text{topology}
\longleftrightarrow
\text{algebra}.
}
\]

%%%%%%%%%%%%%%%%%%%%%%%%%%%%%%%%%%%%%%%%%%%%%%%%%%%%%%%%%%%%
\subsection{Section Connections}
\label{subsec:algebraic-topology-connections}
%%%%%%%%%%%%%%%%%%%%%%%%%%%%%%%%%%%%%%%%%%%%%%%%%%%%%%%%%%%%

\chapterconnection{
Algebraic Topology builds on Point-Set Topology by replacing qualitative
questions about spaces with algebraic invariants. Paths and quotient spaces
lead to homotopy and the fundamental group. Group Theory provides the
algebraic structure of loops, while Linear Algebra and Homological Algebra
provide the machinery for chain complexes and homology. Covering spaces
connect local topology with global structure. Homology and cohomology will
reappear throughout Differential Topology, Geometric Topology, Manifold
Theory, and Knot Theory. The Euler characteristic connects this section with
the earlier study of polyhedra and surfaces. Differential forms later connect
cohomology with Differential Geometry through the de Rham theorem. These
ideas also provide the algebraic foundation for modern applications such as
persistent homology and topological data analysis.
}

%%%%%%%%%%%%%%%%%%%%%%%%%%%%%%%%%%%%%%%%%%%%%%%%%%%%%%%%%%%%
% SECTION SOURCE TRACKING
%%%%%%%%%%%%%%%%%%%%%%%%%%%%%%%%%%%%%%%%%%%%%%%%%%%%%%%%%%%%

%------------------------------------------------------------
% 4.2 Algebraic Topology
%------------------------------------------------------------
%
% Source basis:
%   - Standard introductory algebraic topology,
%     simplicial homology, singular homology, fundamental
%     groups, covering spaces, and cohomology.
%
% Used for:
%   - paths and loops;
%   - path concatenation and reversal;
%   - homotopies;
%   - homotopy classes;
%   - homotopy equivalence;
%   - contractible spaces;
%   - retracts and deformation retracts;
%   - fundamental groups;
%   - winding numbers;
%   - simply connected spaces;
%   - induced homomorphisms;
%   - functoriality;
%   - homotopy invariance;
%   - torus and wedge-of-circles examples;
%   - Seifert--van Kampen theorem;
%   - surface fundamental groups;
%   - covering spaces;
%   - path and homotopy lifting;
%   - universal covers;
%   - deck transformations;
%   - higher homotopy groups;
%   - simplices and simplicial complexes;
%   - chain groups;
%   - boundary operators;
%   - chain complexes;
%   - cycles and boundaries;
%   - simplicial homology;
%   - singular homology;
%   - reduced homology;
%   - coefficient groups;
%   - Betti numbers;
%   - torsion;
%   - boundary matrices;
%   - Euler characteristic;
%   - exact sequences;
%   - relative homology;
%   - Mayer--Vietoris;
%   - excision;
%   - CW complexes;
%   - cochains and cohomology;
%   - cup products;
%   - degree;
%   - Brouwer fixed-point theorem;
%   - no-retraction theorem;
%   - Borsuk--Ulam theorem;
%   - relation between pi_1 and H_1;
%   - persistent homology;
%   - connections with manifolds, knots, and
%     differential forms.
%
% Notes:
%   - Group theory, free groups, presentations,
%     abelianization, exact sequences, and homological
%     algebra are developed canonically in Algebra.
%   - Point-set topology and quotient spaces are developed
%     canonically in Section 4.1.
%   - Differential forms and de Rham theory connect with
%     Differential Geometry and Analysis.
%   - Persistent homology is introduced here only as an
%     application; computational and statistical aspects
%     belong later in Computational Mathematics and
%     Statistics.
%   - Surface topology and manifold classification are
%     developed more fully in later Topology sections.
%   - Knot groups and knot invariants are developed in the
%     Knot Theory section.
%
%%%%%%%%%%%%%%%%%%%%%%%%%%%%%%%%%%%%%%%%%%%%%%%%%%%%%%%%%%%%